\documentclass{amsart}
\usepackage{amsmath,amssymb,amsthm,hyperref}

\newtheorem{theorem}{Theorem}
\newtheorem{lemma}{Lemma}
\newtheorem{proposition}{Proposition}
\newtheorem{corollary}{Corollary}
\theoremstyle{definition}
\newtheorem{definition}{Definition}
\newtheorem{remark}{Remark}

\DeclareMathOperator{\Aut}{Aut}
\DeclareMathOperator{\Hom}{Hom}
\DeclareMathOperator{\Der}{Der}
\DeclareMathOperator{\Spec}{Spec}
\DeclareMathOperator{\GL}{GL}
\DeclareMathOperator{\id}{id}
\newcommand{\Gm}{\mathbb{G}_m}
\newcommand{\g}{\mathfrak{g}}
\newcommand{\fa}{\mathfrak{a}}
\newcommand{\OO}{\mathcal{O}}
\newcommand{\Z}{\mathbb{Z}}
\newcommand{\C}{\mathbb{C}}
\newcommand{\M}{M}

\begin{document}

\title{The automorphism group of an equivariant map algebra}
\author{}
\date{}
\maketitle

\begin{abstract}
We determine the group of Lie algebra automorphisms of an equivariant map
algebra $\M(X,\fa)^\Gamma=(\OO(X)\otimes_k\fa)^\Gamma$ under hypotheses that
hold in all the standard examples (current, loop, multiloop, generalized
current algebras).  The main structural result expresses
$\Aut\bigl(\M(X,\fa)^\Gamma\bigr)$ as a group of \emph{equivariant
twisted automorphisms} of the pair $(X,\fa)$, and as a consequence yields
an explicit, completely computed description for twisted loop algebras and
for multiloop algebras.  We then relate this to the isomorphism
classification of equivariant map algebras.
\end{abstract}

\section{Setup, hypotheses, and statement of results}\label{sec:setup}

Throughout, $k$ is a field of characteristic $0$, $X$ is an affine
$k$-variety (or scheme of finite type) with coordinate ring $A:=\OO(X)$,
and $\fa$ is a Lie algebra over $k$.  A group $\Gamma$ acts on $X$ by
$k$-variety automorphisms (equivalently on $A$ by $k$-algebra
automorphisms) and on $\fa$ by Lie algebra automorphisms.  We write the
$\Gamma$-action on $A$ on the right of the action on $\fa$ via the
\emph{diagonal action} on
\[
\M:=\M(X,\fa)=A\otimes_k\fa ,\qquad
\gamma\cdot(p\otimes u)=(\gamma\cdot p)\otimes(\gamma\cdot u),
\]
which, after the identification $\M=\{$regular maps $X\to\fa\}$, is exactly
$(\gamma\cdot f)(x)=\gamma\cdot f(\gamma^{-1}x)$.  The bracket on $\M$ is
$A$-bilinear:
\[
[p\otimes u,\;q\otimes v]=pq\otimes[u,v].
\]
The equivariant map algebra is the fixed-point Lie subalgebra
\[
\M^\Gamma=\M(X,\fa)^\Gamma=\{f\in\M:\gamma\cdot f=f\ \forall\gamma\in\Gamma\}.
\]

We will work under the following standing hypotheses, which we verify in
the relevant examples in \S\ref{sec:loop}--\S\ref{sec:multiloop}.

\begin{itemize}
\item[(H1)] $A=\OO(X)$ is a finitely generated commutative $k$-algebra and
$X$ is reduced; $k$ is algebraically closed (so that closed points of $X$
are $k$-points and Hilbert's Nullstellensatz applies).
\item[(H2)] $\Gamma$ is finite and acts \emph{freely} on the set $X(k)$ of
$k$-points of $X$.
\item[(H3)] $\fa=\g$ is a finite-dimensional simple Lie algebra over $k$,
and $\Gamma$ acts on $\g$ by automorphisms.
\item[(H4)] Every variety automorphism of $X/\Gamma$ arising from an
automorphism of $\M^\Gamma$ (via Lemma~\ref{lem:centroid}) lifts to a
$k$-variety automorphism $\varphi$ of $X$ normalizing the $\Gamma$-action.
This is needed only to produce the ``variety part'' of an automorphism
(Lemma~\ref{lem:varietypart}); it is automatic in every example treated
here (loop, multiloop: the cover is a power map and the relevant
$\bar\varphi$ are monomial) and holds whenever $X/\Gamma$ is a torus, by
Pianzola's $H^1$-vanishing for Laurent rings.
\end{itemize}

Because $\Gamma$ is finite and $\mathrm{char}\,k=0$, the functor of taking
$\Gamma$-invariants is exact and commutes with $\otimes_k$ on
$\Gamma$-modules; the averaging idempotent
$e=\frac1{|\Gamma|}\sum_{\gamma\in\Gamma}\gamma$ is the projector onto
invariants.

\medskip
We now introduce the group that will turn out to be the automorphism
group.

\begin{definition}\label{def:emap}
Let $G:=\Aut_{\mathrm{Lie}}(\g)$ be the (linear algebraic) automorphism
group of $\g$, with subgroup of inner automorphisms $G^\circ=\mathrm{Inn}(\g)$
(the identity component when $k=\C$), and let
$\mathrm{Out}(\g)=G/G^\circ$ be the (finite) group of outer
automorphisms.  We write $\gamma$ also for the image of $\gamma\in\Gamma$ in
$G=\Aut(\g)$ (the subgroup $\Gamma_\g\subseteq G$); and we set
\[
N_\Gamma(\g):=\{\,\theta\in G:
\theta\,\Gamma_\g\,\theta^{-1}=\Gamma_\g\,\},
\]
the normalizer in $G$ of $\Gamma_\g$.

A \emph{$\Gamma$-equivariant twisted automorphism} of $(X,\g)$ is a pair
$(\varphi,\Psi)$ where $\varphi\colon X\to X$ is a $k$-variety
automorphism normalizing the $\Gamma$-action (i.e.\ $\varphi\circ\gamma=
c(\gamma)\circ\varphi$ for some $c(\gamma)\in\Gamma$ for each
$\gamma$, where $c\in\Aut(\Gamma)$ is a fixed automorphism) and
$\Psi\colon X\to G$ is a regular map (a morphism of varieties to the
algebraic group $G$) satisfying the \emph{twisted equivariance}
\begin{equation}\label{eq:equivdef}
\Psi\bigl(c(\gamma)x\bigr)=c(\gamma)\,\Psi(x)\,\gamma^{-1}
\qquad(\gamma\in\Gamma,\ x\in X(k)),
\end{equation}
where on the right $c(\gamma),\gamma\in G$ via $\Gamma_\g$.  The associated
operator is
\begin{equation}\label{eq:Phi}
\bigl(\Phi f\bigr)(x):=\Psi(x)\,\bigl(f(\varphi^{-1}x)\bigr),
\qquad f\in\M,\ x\in X(k).
\end{equation}
We call $\Phi$ the operator \emph{induced} by $(\varphi,\Psi)$, and we let
$\mathcal{E}(X,\g)^\Gamma$ denote the group of such operators on
$\M^\Gamma$.

\medskip
\noindent\textbf{Important correction.}  The gauge map $\Psi$ is valued in
the \emph{full} automorphism group $G=\Aut(\g)$, \emph{not} merely in
$N_\Gamma(\g)$.  Indeed \eqref{eq:equivdef} only relates the values of
$\Psi$ at the \emph{distinct} points of a $\Gamma$-orbit: on a free orbit
$\{x_1,\dots,x_{|\Gamma|}\}$ with $c=\id$ it reads
$\Psi(\gamma x_1)=\gamma\,\Psi(x_1)\,\gamma^{-1}$, so $\Psi(x_1)\in G$ may
be \emph{arbitrary} and the remaining values are then determined.  (For
example with $\g=\mathfrak{sl}_2$ and $\Gamma=\Z/m$ acting by a nontrivial
$\sigma$, taking $\Psi(x_1)=\theta$ a generic inner automorphism not
normalizing $\langle\sigma\rangle$ still yields a well-defined Lie
automorphism of $\M^\Gamma$.)  Thus the correct gauge group is the group of
twisted-equivariant $G$-valued maps
\begin{equation}\label{eq:gaugegroup}
\mathcal G_\Gamma(X,\g)\ :=\ \bigl\{\Psi\colon X\to G\ \text{regular}\ :\
\eqref{eq:equivdef}\ \text{holds}\bigr\},
\end{equation}
a group under pointwise multiplication.  The subgroup $N_\Gamma(\g)$ of
\emph{constant} admissible gauges appears only as part of the component
group; see Remark~\ref{rem:groupstr}.
\end{definition}

\begin{theorem}[Main structural theorem]\label{thm:main}
Assume \textnormal{(H1)--(H4)}.  Then every $\Phi\in\mathcal{E}(X,\g)^\Gamma$
is a Lie algebra automorphism of $\M^\Gamma$, and
\[
\Aut_{\mathrm{Lie}}\bigl(\M(X,\g)^\Gamma\bigr)
=\mathcal{E}(X,\g)^\Gamma .
\]
Equivalently: every Lie algebra automorphism of $\M^\Gamma$ is induced,
via \eqref{eq:Phi}, by a unique pair $(\varphi,\Psi)$ consisting of a
$k$-variety automorphism $\varphi$ of $X$ normalizing $\Gamma$ and a
regular map $\Psi\colon X\to G=\Aut(\g)$ subject to the twisted
equivariance \eqref{eq:equiv}.  In group-theoretic terms (Remark~\ref{rem:groupstr}),
\[
\Aut_{\mathrm{Lie}}\bigl(\M(X,\g)^\Gamma\bigr)\;\cong\;
\Aut(X,\Gamma)\ \ltimes\ \mathcal G_\Gamma(X,\g),
\]
the semidirect product of the variety automorphisms normalizing $\Gamma$
with the twisted equivariant gauge group $\mathcal G_\Gamma(X,\g)$ of
\eqref{eq:gaugegroup}.
\end{theorem}

The proof occupies \S\ref{sec:central}--\S\ref{sec:autthm}.  The
explicit computations are:

\begin{theorem}[Twisted loop algebras]\label{thm:loop}
Let $X=\Gm$, let $\g$ be finite-dimensional simple over an algebraically
closed $k$, let $\sigma$ be a diagram automorphism of order $m$ and
$\Gamma=\Z/m\Z=\langle\gamma\rangle$ act by $\gamma\cdot=\sigma$ on $\g$ and
by $t\mapsto\zeta^{-1}t$ on $\Gm$ (with $\zeta$ a fixed primitive $m$-th root
of unity).  Write $L:=\M(\Gm,\g)^\Gamma=L(\g,\sigma)$ for the twisted loop
algebra.  Then, as a special case of Theorem~\ref{thm:main},
\[
\Aut_{\mathrm{Lie}}(L)\;\cong\;\mathcal E(\Gm,\g)^\Gamma
\;\cong\;
\Aut(\Gm,\Gamma)\ \ltimes\ \mathcal G_\Gamma(\Gm,\g),
\]
where (see \S\ref{sec:loop} for the verifications):
\begin{enumerate}
\item the variety part is
$\Aut(\Gm,\Gamma)=k^\times\rtimes\Z/2$, with $k^\times$ the rescalings
$\varphi_a\colon t\mapsto at$ (acting by $c=\id$ on $\Gamma$) and the
generator of $\Z/2$ the inversion $\iota\colon t\mapsto t^{-1}$ (acting by
$c(\gamma)=\gamma^{-1}$, hence forcing $\sigma\mapsto\sigma^{-1}$ on
$\g$, consistently with $L(\g,\sigma)\cong L(\g,\sigma^{-1})$);
\item the gauge part $\mathcal G_\Gamma(\Gm,\g)$ consists of
the regular (Laurent-polynomial) maps $\Psi\colon\Gm\to G=\Aut(\g)$
satisfying \eqref{eq:equiv}; concretely, for the $c=\id$ component this is
$\Psi(\zeta x)=\sigma^{-1}\Psi(x)\sigma$.  This is the \emph{twisted loop
group} $G(\sigma):=\{\Psi\in G(k[t^{\pm1}]):\Psi(\zeta t)=
\sigma^{-1}\Psi(t)\sigma\}$ attached to $(G,\sigma)$; it is \emph{not}
$N_\Gamma(\g)$-valued (only its constant part is).
\end{enumerate}
Consequently every automorphism of $L$ has the explicit form
\begin{equation}\label{eq:loopaut}
(\Phi f)(t)=\Psi(t)\bigl(f(a\,t^{\pm1})\bigr),\qquad a\in k^\times,
\end{equation}
and the group is generated by:
\begin{enumerate}
\item the constant gauge maps that are admissible together with some
variety part: for $\varphi$ with $c=\id$ a constant $\Psi\equiv\theta$
requires $\sigma\theta\sigma^{-1}=\theta$, i.e.\ $\theta\in Z_G(\sigma)$;
allowing $\varphi$ to induce $c(\gamma)=\gamma^{r}$ (a power), the
admissible constants form $N_\Gamma(\g)=\Aut(\g,\langle\sigma\rangle):=
\{\theta\in\Aut(\g):\theta\langle\sigma\rangle\theta^{-1}=
\langle\sigma\rangle\}$, the diagram/outer symmetries of the pair
$(\g,\langle\sigma\rangle)$ (see \S\ref{sec:loop});
\item the rescalings $\varphi_a\colon t\mapsto a t$, $a\in k^\times$;
\item the inversion $\iota\colon t\mapsto t^{-1}$, accompanied by
$\sigma\mapsto\sigma^{-1}$ on $\g$;
\item the equivariant inner gauge loops, which together generate the inner
automorphism group $\mathrm{Inn}(L)$.
\end{enumerate}
\end{theorem}

\begin{theorem}[Multiloop algebras]\label{thm:multiloop}
Let $X=(\Gm)^n$, $\Gamma=\Z/m_1\times\cdots\times\Z/m_n$ acting on $X$ by
$\gamma_j\colon t_i\mapsto \zeta_{m_j}^{-\delta_{ij}}t_i$ and on $\g$ by
commuting finite-order automorphisms $\sigma_1,\dots,\sigma_n$.  Let
$M=M(\g,\sigma_1,\dots,\sigma_n)=\M(X,\g)^\Gamma$ be the corresponding
multiloop algebra.  Then every automorphism of $M$ is induced by a pair
$(\varphi,\Psi)$ with $\varphi$ a monomial automorphism of $(\Gm)^n$
normalizing $\Gamma$ and $\Psi\colon(\Gm)^n\to G=\Aut(\g)$ a regular gauge
loop, subject to the twisted equivariance \eqref{eq:equiv}; as in the loop
case $\Psi$ is $G$-valued, not merely $N_\Gamma(\g)$-valued.  Concretely
\[
\Aut_{\mathrm{Lie}}(M)\cong
\bigl((k^\times)^n\rtimes\GL_n(\Z)\bigr)_\Gamma
\ \ltimes\ \mathcal G_\Gamma\bigl((\Gm)^n,\g\bigr),
\]
and the gauge group decomposes as $\mathcal G_\Gamma=\mathrm{Inn}(M)\cdot R$
with $R$ the finite outer/diagram part coming from constant gauges in
$N_\Gamma(\g)$, where
$N_\Gamma(\g)=\Aut(\g;\sigma_1,\dots,\sigma_n):=\{\theta\in\Aut(\g):
\theta\langle\sigma_1,\dots,\sigma_n\rangle\theta^{-1}=
\langle\sigma_1,\dots,\sigma_n\rangle\}$ is the group of $\theta\in\Aut(\g)$
permuting and powering the $\sigma_i$ compatibly with a monomial $\varphi$
(the matrix $B\in\GL_n(\Z)$ realizing the same action on the character
lattice $\Gamma^\vee=\prod_j\Z/m_j$).
\end{theorem}

\section{Centroid, evaluations, and ideals}\label{sec:central}

The proof of Theorem~\ref{thm:main} rests on three rigidity facts:
the centroid of $\M^\Gamma$ recovers $A^\Gamma$ (whence the variety part
$\varphi$), the maximal ideals of finite codimension recover the points of
$X/\Gamma$, and a fibrewise simplicity statement lets us reconstruct
$\Psi$.  We prove them now.

\subsection{The centroid}
For a Lie algebra $\mathfrak l$, its \emph{centroid} is
\[
C(\mathfrak l)=\{c\in\End_k(\mathfrak l):
c[x,y]=[cx,y]=[x,cy]\ \forall x,y\}.
\]
$C(\mathfrak l)$ is a commutative $k$-algebra when $[\mathfrak l,\mathfrak
l]=\mathfrak l$, and any Lie automorphism $\Phi$ of $\mathfrak l$ induces a
$k$-algebra automorphism $c\mapsto\Phi c\Phi^{-1}$ of $C(\mathfrak l)$.

\begin{lemma}\label{lem:perfect}
Under \textnormal{(H1)--(H3)}, $\M^\Gamma$ is perfect, i.e.\
$[\M^\Gamma,\M^\Gamma]=\M^\Gamma$.
\end{lemma}
\begin{proof}
Since $\g$ is simple, $[\g,\g]=\g$.  First, $\M=A\otimes\g$ is perfect: for
$p\otimes u\in\M$ write $u=\sum_i[v_i,w_i]$ (possible as $[\g,\g]=\g$); then
$p\otimes u=\sum_i[1\otimes v_i,\,p\otimes w_i]\in[\M,\M]$, so $[\M,\M]=\M$.

Now $[\M,\M]=\M$ is an equality of $\Gamma$-submodules of $\M$ (the bracket
is $\Gamma$-equivariant, so $[\M,\M]$ is a $\Gamma$-submodule).  Apply the
exact functor $(-)^\Gamma$ (exact because $\Gamma$ is finite and
$\mathrm{char}\,k=0$, with projector $e$): we get
\begin{equation}\label{eq:perf1}
\M^\Gamma=[\M,\M]^\Gamma .
\end{equation}
It remains to prove $[\M,\M]^\Gamma=[\M^\Gamma,\M^\Gamma]$.  The inclusion
$\supseteq$ is immediate since $\M^\Gamma$ is a subalgebra.  For
$\subseteq$, let $z\in[\M,\M]^\Gamma$; then $z=ez$ where, writing
$z=\sum_i[x_i,y_i]$,
\[
z=ez=\tfrac1{|\Gamma|}\sum_{\gamma\in\Gamma}\gamma\!\cdot\! z
=\tfrac1{|\Gamma|}\sum_{\gamma}\sum_i[\gamma\!\cdot\!x_i,\;
\gamma\!\cdot\!y_i],
\]
using $\Gamma$-equivariance of the bracket.  This realizes $z$ as a sum of
brackets, but not yet of \emph{invariant} elements.  To fix this, use the
$A^\Gamma$-bilinearity of the bracket and the fact (H2) that $A$ is a
finitely generated projective $A^\Gamma$-module: choose a finite
\emph{averaging dual basis}, i.e.\ elements $p_\alpha,q_\alpha\in A$ with
\begin{equation}\label{eq:dualbasis}
\sum_\alpha p_\alpha\,\gamma(q_\alpha)=\delta_{\gamma,1}\in A
\qquad(\gamma\in\Gamma),
\end{equation}
where $\delta_{\gamma,1}$ is the constant function $1$ if $\gamma=1$ and
$0$ otherwise (such a system exists: by (H2) the quotient $X\to X/\Gamma$
is a finite \'etale Galois cover with group $\Gamma$, hence $A$ is a
$\Gamma$-Galois extension of $A^\Gamma$ and admits a normal basis, which
yields a self-dual averaging system \eqref{eq:dualbasis} over $A^\Gamma$).
Then for any $a,b\in A$ and $u,v\in\g$,
\begin{equation}\label{eq:brkid}
e\bigl[a\otimes u,\,b\otimes v\bigr]
=|\Gamma|\sum_\alpha\Bigl[\,e\bigl((a p_\alpha)\otimes u\bigr),\;
e\bigl((b q_\alpha)\otimes v\bigr)\Bigr],
\end{equation}
a finite sum of brackets of \emph{invariant} elements
$e((ap_\alpha)\otimes u),\,e((bq_\alpha)\otimes v)\in\M^\Gamma$.  To
verify \eqref{eq:brkid}, expand the right side using $A^\Gamma$-bilinearity
of the bracket and $e(c\otimes w)=\frac1{|\Gamma|}\sum_{\delta}
\delta(c)\otimes\delta(w)$:
\[
|\Gamma|\sum_\alpha\bigl[e((ap_\alpha)\otimes u),e((bq_\alpha)\otimes v)\bigr]
=\frac1{|\Gamma|}\sum_{\delta,\delta'}\sum_\alpha
\bigl[\delta(ap_\alpha)\otimes\delta(u),\,
\delta'(bq_\alpha)\otimes\delta'(v)\bigr].
\]
For fixed $\delta,\delta'$ the inner sum over $\alpha$ contributes the
$A$-coefficient $\sum_\alpha\delta(ap_\alpha)\delta'(bq_\alpha)
=\delta(a)\delta'(b)\sum_\alpha\delta(p_\alpha)\delta'(q_\alpha)
=\delta(a)\delta'(b)\,\delta\!\Bigl(\sum_\alpha p_\alpha\,
(\delta^{-1}\delta')(q_\alpha)\Bigr)
=\delta(a)\delta'(b)\,\delta(\delta_{\delta^{-1}\delta',1})$, which by
\eqref{eq:dualbasis} vanishes unless $\delta'=\delta$ and equals
$\delta(ab)$ when $\delta'=\delta$.  Hence only the diagonal terms
$\delta'=\delta$ survive and the right side becomes
$\frac1{|\Gamma|}\sum_\delta[\delta(ab)\otimes\delta(u),\,1\otimes
\delta(v)]=\frac1{|\Gamma|}\sum_\delta\delta(ab)\otimes[\delta(u),\delta(v)]
=\frac1{|\Gamma|}\sum_\delta\delta\bigl(ab\otimes[u,v]\bigr)
=e[a\otimes u,b\otimes v]$, proving \eqref{eq:brkid}.  Applying
\eqref{eq:brkid} to each summand of $z=\sum_i[x_i,y_i]$ (after writing each
$x_i,y_i$ as a sum of simple tensors $a\otimes u$) shows
$z=ez\in[\M^\Gamma,\M^\Gamma]$.  With \eqref{eq:perf1} this gives
$\M^\Gamma=[\M^\Gamma,\M^\Gamma]$.
\end{proof}

\begin{lemma}\label{lem:centroid}
Under \textnormal{(H1)--(H3)}, the natural map
\[
A^\Gamma=\OO(X)^\Gamma=\OO(X/\Gamma)\;\longrightarrow\;C(\M^\Gamma),
\qquad p\mapsto(f\mapsto pf),
\]
is an isomorphism of $k$-algebras.  Consequently every Lie automorphism
$\Phi$ of $\M^\Gamma$ induces a $k$-algebra automorphism $\varphi^*$ of
$\OO(X/\Gamma)$, i.e.\ a variety automorphism $\bar\varphi$ of $X/\Gamma$.
\end{lemma}
\begin{proof}
Multiplication by $p\in A^\Gamma$ is $A$-linear hence centroidal, and is
$\Gamma$-equivariant, so it preserves $\M^\Gamma$ and lies in
$C(\M^\Gamma)$; the map is injective because $\M^\Gamma$ is a faithful
$A^\Gamma$-module: if $p\in A^\Gamma$, $p\ne0$, then by the base-change
isomorphism \eqref{eq:basechange} (proved below) $A\otimes_{A^\Gamma}
\M^\Gamma\cong\M=A\otimes_k\g$, and multiplication by $p$ on $\M$ is
nonzero (as $A$ is reduced and $\g\ne0$); since $A$ is faithfully flat over
$A^\Gamma$, multiplication by $p$ on $\M^\Gamma$ is already nonzero.
Hence $p\mapsto(f\mapsto pf)$ is injective.  For surjectivity: $\g$ is central-simple, so its
centroid is $k$; the centroid of $A\otimes_k\g$ is then
$C(A\otimes\g)=A\otimes C(\g)=A$ by the standard computation
(any $c\in C(A\otimes\g)$ commutes with $\mathrm{ad}$ of all $1\otimes u$,
and using $[\g,\g]=\g$ together with simplicity forces $c=$ multiplication
by some $p\in A$).  We now pass to invariants.  By (H2) the geometric quotient
$\pi\colon X\to X/\Gamma$ is a finite \'etale Galois cover with group
$\Gamma$ (a finite group acting freely on a smooth-enough affine variety in
characteristic $0$ gives an \'etale quotient; freeness on $k$-points and
reducedness (H1) give the \'etale, hence flat, property).  Consequently
$A$ is a faithfully flat, finite projective $A^\Gamma$-module of constant
rank $|\Gamma|$, and the multiplication map induces a canonical
isomorphism of Lie algebras
\begin{equation}\label{eq:basechange}
A\otimes_{A^\Gamma}\M^\Gamma\;\xrightarrow{\ \sim\ }\;\M ,
\qquad p\otimes f\mapsto p\,f .
\end{equation}
Indeed, the map is $A$-linear, and it suffices to check it is an
isomorphism after the faithfully flat base change $A^\Gamma\to A$.  Galois
descent for the $\Gamma$-Galois extension $A^\Gamma\subseteq A$ gives the
$\Gamma$-equivariant isomorphism $A\otimes_{A^\Gamma}A\cong
\prod_{\gamma\in\Gamma}A$ (the normal-basis / Galois isomorphism).  Since
$\Gamma$ is finite and $\operatorname{char}k=0$, taking $\Gamma$-invariants
commutes with the flat base change $A^\Gamma\to A$, so
\[
A\otimes_{A^\Gamma}\M^\Gamma
=A\otimes_{A^\Gamma}(A\otimes_k\g)^\Gamma
\cong\bigl(A\otimes_{A^\Gamma}A\otimes_k\g\bigr)^\Gamma
\cong\Bigl(\textstyle\prod_{\gamma}A\otimes_k\g\Bigr)^\Gamma
\cong A\otimes_k\g=\M,
\]
the last isomorphism being projection to the diagonal factor (as in
Lemma~\ref{lem:fibres}, $\Gamma$ permutes the $|\Gamma|$ factors simply
transitively, so the invariants are the ``diagonal'' copy $A\otimes_k\g$).
Tracing through, this composite is exactly multiplication
$p\otimes f\mapsto pf$; this proves \eqref{eq:basechange}.

We claim $C(\M^\Gamma)=A^\Gamma$.  The inclusion
$A^\Gamma\subseteq C(\M^\Gamma)$ and injectivity were shown above.  For the
reverse inclusion, let $c\in C(\M^\Gamma)$.  Because $\M^\Gamma$ is perfect
(Lemma~\ref{lem:perfect}), $c$ is $A^\Gamma$-linear: for $p\in A^\Gamma$
and $x=\sum_i[y_i,z_i]\in\M^\Gamma$ we have
$c(px)=\sum_i c([py_i,z_i])=\sum_i[c(py_i),z_i]$ and also
$pc(x)=\sum_i p[c(y_i),z_i]=\sum_i[pc(y_i),z_i]=\sum_i[c(py_i),z_i]$ since
$p$ is centroidal; hence $c(px)=pc(x)$.  Therefore
$c\otimes\id_A\colon A\otimes_{A^\Gamma}\M^\Gamma\to
A\otimes_{A^\Gamma}\M^\Gamma$ is defined, and under \eqref{eq:basechange}
it is an $A$-linear centroidal operator $\tilde c$ on $\M=A\otimes_k\g$
which restricts to $c$ on $\M^\Gamma$ and is $\Gamma$-equivariant (it is
the base change of an operator defined over $A^\Gamma$).  By the
computation of the previous paragraph, $C(\M)=A$, so $\tilde c$ is
multiplication by some $\tilde p\in A$; its $\Gamma$-equivariance forces
$\tilde p\in A^\Gamma$, and restricting to $\M^\Gamma$ gives
$c=\tilde p\cdot$ with $\tilde p\in A^\Gamma$.  Hence
$C(\M^\Gamma)=A^\Gamma$.
\end{proof}

\subsection{Evaluation ideals and the variety part}
For a $k$-point $x\in X(k)$ let $\mathfrak m_x\subset A$ be its maximal
ideal and $\mathrm{ev}_x\colon\M\to\g$, $f\mapsto f(x)$.  For
$\bar x\in (X/\Gamma)(k)$ let $\mathfrak m_{\bar x}\subset A^\Gamma$ be its
maximal ideal and $I_{\bar x}=\mathfrak m_{\bar x}\,\M^\Gamma$.

\begin{lemma}\label{lem:fibres}
Assume \textnormal{(H1)--(H3)}.  For each $\bar x\in (X/\Gamma)(k)$ with
$\Gamma$-orbit $\{x_1,\dots,x_{|\Gamma|}\}\subset X(k)$:
\begin{enumerate}
\item The quotient $\M^\Gamma/I_{\bar x}$ is, as a Lie algebra,
canonically isomorphic to $\g$ via $\mathrm{ev}_{x_1}$ (restricted to
$\M^\Gamma$); in particular it is simple.
\item $I_{\bar x}$ is a maximal ideal of $\M^\Gamma$ of finite
codimension $\dim_k\g$, and these exhaust the finite-codimension maximal
ideals with simple semisimple quotient.
\end{enumerate}
\end{lemma}
\begin{proof}
Because the $\Gamma$-action on $X(k)$ is free (H2), the orbit of $x_1$ has
exactly $|\Gamma|$ distinct points, the local rings $A_{x_i}$ are permuted
freely, and by the Chinese Remainder Theorem the restriction map
\[
\M^\Gamma\longrightarrow\Bigl(\textstyle\prod_{i}\g\Bigr)^\Gamma\cong\g,
\qquad f\mapsto f(x_1),
\]
is surjective with kernel $I_{\bar x}$.  We first record the precise form
of the target.  Since the orbit $\{x_1,\dots,x_{|\Gamma|}\}$ consists of
$|\Gamma|$ distinct $k$-points (freeness, (H2)), the evaluation map
$\mathrm{ev}\colon\M=A\otimes\g\to\prod_{i}\g$, $f\mapsto(f(x_i))_i$, is
$\Gamma$-equivariant for the diagonal action on the source and the action
$\gamma\cdot(w_i)_i=(\gamma\cdot w_{\gamma^{-1}i})_i$ on
$\prod_i\g$, where $\gamma$ permutes the orbit by
$x_{\gamma i}=\gamma x_i$.  Taking $\Gamma$-invariants and using that
$\Gamma$ permutes the $|\Gamma|$ factors \emph{simply transitively}, the
invariants $\bigl(\prod_i\g\bigr)^\Gamma$ are exactly the tuples
determined by their value at $x_1$: the projection
$\mathrm{pr}_{1}\colon\bigl(\prod_i\g\bigr)^\Gamma\to\g$ is an
isomorphism (its inverse sends $u\in\g$ to the tuple with
$w_{\gamma 1}=\gamma\cdot u$, well defined precisely because the
stabiliser of $x_1$ is trivial).  Thus it suffices to show
$\mathrm{ev}|_{\M^\Gamma}\colon\M^\Gamma\to\bigl(\prod_i\g\bigr)^\Gamma$
is onto; equivalently, that $\mathrm{ev}_{x_1}|_{\M^\Gamma}$ is onto $\g$.
Given $u\in\g$, choose $p\in A$ with $p(x_1)=1$ and $p(x_i)=0$ for
$i\ne1$ (possible by the Chinese Remainder Theorem / Nullstellensatz,
(H1), as the $\mathfrak m_{x_i}$ are pairwise comaximal), and set
$f=e(p\otimes u)=\frac1{|\Gamma|}\sum_{\gamma}\gamma\cdot(p\otimes u)
\in\M^\Gamma$.  Then
\[
f(x_1)=\frac1{|\Gamma|}\sum_\gamma (\gamma\cdot p)(x_1)\,\gamma\cdot u
=\frac1{|\Gamma|}\sum_\gamma p(\gamma^{-1}x_1)\,\gamma\cdot u .
\]
Because the orbit is free, $p(\gamma^{-1}x_1)$ equals
$1$ for $\gamma=1$ and $0$ otherwise, so $f(x_1)=\frac1{|\Gamma|}u$.
Replacing $u$ by $|\Gamma|u$ shows $\mathrm{ev}_{x_1}|_{\M^\Gamma}$ is
surjective, hence so is $\mathrm{ev}|_{\M^\Gamma}$ onto the invariants.  The kernel is exactly
$\{f:f(x_1)=0\}=\{f:f|_{\text{orbit}}=0\}=\mathfrak m_{\bar x}\M^\Gamma$,
the last equality because $\mathfrak m_{\bar x}A=\bigcap_i\mathfrak m_{x_i}$
(\'etale descent / freeness).  Since the quotient is the simple Lie
algebra $\g$, $I_{\bar x}$ is maximal of codimension $\dim\g$.

\emph{These exhaust the relevant ideals.}  Let $I\subseteq\M^\Gamma$ be a
Lie ideal of finite codimension with $\M^\Gamma/I$ a (nonzero) simple Lie
algebra.  We first promote $I$ to an ideal of $\M$.  Set
$\widehat I:=A\cdot I\subseteq\M$, the $A$-submodule of $\M$ generated by
$I$ (using the $A$-module structure on $\M=A\otimes_k\g$ and the inclusion
$\M^\Gamma\subseteq\M$).  Because $A$ is flat over $A^\Gamma$ and
$\M\cong A\otimes_{A^\Gamma}\M^\Gamma$ by the base-change isomorphism
\eqref{eq:basechange}, applying $A\otimes_{A^\Gamma}(-)$ to the exact
sequence $0\to I\to\M^\Gamma\to\M^\Gamma/I\to0$ of $A^\Gamma$-modules gives
the exact sequence of $A$-modules
\begin{equation}\label{eq:idealbasechange}
0\to\widehat I\to\M\to A\otimes_{A^\Gamma}(\M^\Gamma/I)\to0 ,
\qquad\widehat I\cong A\otimes_{A^\Gamma}I .
\end{equation}
Now $\widehat I=A\cdot I$ is a Lie ideal of $\M$: the bracket of $\M$ is
$A$-bilinear, so for $a\in A$, $y\in I$ and $z\in\M$ we have
$[z,\,a\,y]=a\,[z,\,y]$.  Writing $z=\sum_k b_k z_k$ with $z_k\in\M^\Gamma$
($A$-spanning by \eqref{eq:basechange}) and using that $I$ is an ideal of
$\M^\Gamma$ gives $[z_k,y]\in I$, whence $[z,ay]=\sum_k a b_k[z_k,y]\in
A\cdot I=\widehat I$.  Thus $\widehat I$ absorbs brackets.  By
\eqref{eq:idealbasechange} the quotient $\M/\widehat I\cong
A\otimes_{A^\Gamma}(\M^\Gamma/I)$ is finite-dimensional over $k$, since
$\M^\Gamma/I$ is finite-dimensional and $A$ is a finite $A^\Gamma$-module
(rank $|\Gamma|$, (H2)).  Thus $\widehat I$ is a finite-codimension Lie
ideal of $\M=A\otimes_k\g$.

We now classify the finite-codimension Lie ideals of $\M=A\otimes_k\g$
($\g$ simple).  Let $J\subseteq A\otimes\g$ be such an ideal and put
$\mathfrak a_J:=\{c\in A:c\otimes\g\subseteq J\}$.  We claim $\mathfrak
a_J$ is an ideal of $A$ and $J=\mathfrak a_J\otimes\g$.  First, if
$b\otimes u'\in J$ with $u'\ne0$ then for every $a\in A$ and every
$u''\in\g$, $A$-bilinearity gives $ab\otimes[u'',u']=[a\otimes u'',\,
b\otimes u']\in J$; as $\g$ is simple and $u'\ne0$ we have $[\g,u']=\g$, so
$ab\otimes\g\subseteq J$ for all $a\in A$.  In particular
$b\in\mathfrak a_J$.  Consequently, fixing a $k$-basis $(u_i)$ of $\g$ and
writing any $w\in J$ as $w=\sum_i c_i\otimes u_i$, each coefficient $c_i$
(obtained by an $A$-linear projection $\M\to A$, which carries $J$ into the
set of coefficients arising from elements of $J$) lies in $\mathfrak a_J$
by the above; hence $J\subseteq\mathfrak a_J\otimes\g$, and the reverse
inclusion is the definition of $\mathfrak a_J$.  Thus $J=\mathfrak a_J
\otimes\g$, and $\mathfrak a_J$ is an ideal of $A$ (closed under
multiplication by $A$ by the displayed computation).  Finite codimension of
$J=\mathfrak a_J\otimes\g$ forces $\dim_k(A/\mathfrak a_J)<\infty$.  In
particular $A/\mathfrak a_J$ is an Artinian $k$-algebra, hence a finite
product of local Artinian rings; its reduction is a product of copies of
$k$ (algebraically closed, (H1)), corresponding to finitely many points
$x_1,\dots,x_r\in X(k)$ with $\sqrt{\mathfrak a_J}=\bigcap_j\mathfrak
m_{x_j}$.

If in addition $\M/J=(A/\mathfrak a_J)\otimes\g$ is \emph{simple}, then
$A/\mathfrak a_J$ has no proper nonzero ideal $\bar{\mathfrak b}$ (else
$\bar{\mathfrak b}\otimes\g$ would be a proper nonzero ideal of the
quotient), so $A/\mathfrak a_J$ is a field, hence $=k$; thus
$\mathfrak a_J=\mathfrak m_x$ and $J=\mathfrak m_x\otimes\g$ for a unique
$x\in X(k)$.

Apply the ideal classification to $J=\widehat I$.  Since $I\subseteq\M^\Gamma$,
the generated ideal $\widehat I=A\cdot I$ is $\Gamma$-stable inside $\M$ (it is
the image of the $\Gamma$-equivariant base change $A\otimes_{A^\Gamma}I$ under
\eqref{eq:basechange}, $\Gamma$ acting only on the $A$-factor).  By the
classification proved above, $\widehat I=\mathfrak a\otimes\g$ with
$\mathfrak a:=\mathfrak a_{\widehat I}\subseteq A$ an ideal of finite
codimension, and $\mathfrak a$ is $\Gamma$-stable because $\widehat I$ is.  We
claim $\mathfrak a$ is radical and is the ideal of a single $\Gamma$-orbit.

First, $\widehat I\cap\M^\Gamma=I$: by faithful flatness of $A$ over
$A^\Gamma$ and the base-change isomorphism \eqref{eq:idealbasechange}, applying
$(-)^\Gamma$ (exact, projector $e$) to $\widehat I\cong A\otimes_{A^\Gamma}I$
gives $(\widehat I)^\Gamma=(A\otimes_{A^\Gamma}I)^\Gamma\supseteq I$, and the
reverse inclusion $\widehat I\cap\M^\Gamma\subseteq I$ holds because $e$ maps
$\widehat I=A\cdot I$ into $A^\Gamma\cdot I=I$ (using
\eqref{eq:brkid}-style averaging of $A$-multiples of elements of $I$, or simply
$e(\widehat I)=\widehat I^\Gamma$ and faithful flatness).  Hence
\begin{equation}\label{eq:Irecover}
I=\widehat I\cap\M^\Gamma=(\mathfrak a\otimes\g)^\Gamma
=(\mathfrak a^\Gamma\!\cdot A\ \text{part})\ \Longrightarrow\
\M^\Gamma/I\cong(\,\M/\widehat I\,)^\Gamma=\bigl((A/\mathfrak a)\otimes\g\bigr)^\Gamma,
\end{equation}
the last isomorphism again by exactness of $(-)^\Gamma$ applied to
$0\to\widehat I\to\M\to(A/\mathfrak a)\otimes\g\to0$.

Now write the Artinian $\Gamma$-algebra $\bar A:=A/\mathfrak a$ as a product
$\bar A\cong\prod_{x\in S}\bar A_x$ of its local factors, indexed by the finite
$\Gamma$-stable set $S=\{x_1,\dots,x_r\}\subseteq X(k)$ of points in its support
($\sqrt{\mathfrak a}=\bigcap_{x\in S}\mathfrak m_x$); $\Gamma$ permutes the
factors according to its action on $S$.  Then
$\M^\Gamma/I\cong\bigl(\prod_{x\in S}\bar A_x\otimes\g\bigr)^\Gamma$.
By hypothesis this is a \emph{simple} Lie algebra.

We rule out nilpotents and extra orbits.  Suppose the nilradical
$\mathfrak n=\sqrt{\mathfrak a}/\mathfrak a\subseteq\bar A$ were nonzero; being a
$\Gamma$-stable nilpotent ideal of $\bar A$, the subspace
$(\mathfrak n\otimes\g)^\Gamma$ is a nonzero \emph{abelian} (indeed nilpotent)
ideal of $\M^\Gamma/I$ — nonzero because $A$ is faithfully flat over $A^\Gamma$
so taking invariants does not kill the nonzero $\Gamma$-module
$\mathfrak n\otimes\g$.  A simple Lie algebra has no nonzero abelian ideal, a
contradiction; hence $\mathfrak n=0$, i.e.\ $\bar A$ is reduced and
$\bar A\cong\prod_{x\in S}k$.  Thus
$\M^\Gamma/I\cong\bigl(\prod_{x\in S}\g\bigr)^\Gamma$.  If $S$ decomposes
into $\ell$ distinct $\Gamma$-orbits $O_1,\dots,O_\ell$, then
$\bigl(\prod_{x\in S}\g\bigr)^\Gamma\cong\bigoplus_{j=1}^\ell
\bigl(\prod_{x\in O_j}\g\bigr)^\Gamma\cong\g^{\oplus\ell}$
(each orbit block is $\cong\g$ by the freeness computation of
Lemma~\ref{lem:fibres}(1)), which is simple iff $\ell=1$.  Therefore $S$ is a
single free $\Gamma$-orbit $\Gamma x=\{x_1,\dots,x_{|\Gamma|}\}$, and
$\mathfrak a=\bigcap_i\mathfrak m_{x_i}=\mathfrak m_{\bar x}A$ for
$\bar x=\pi(x)$.  Substituting back,
\[
I=\widehat I\cap\M^\Gamma=(\mathfrak m_{\bar x}A\otimes\g)^\Gamma
=\mathfrak m_{\bar x}\,(A\otimes\g)^\Gamma=\mathfrak m_{\bar x}\M^\Gamma=
I_{\bar x},
\]
the third equality because $\mathfrak m_{\bar x}\subseteq A^\Gamma$ is
centroidal and $(-)^\Gamma$ is $A^\Gamma$-linear.
Hence every finite-codimension ideal with simple quotient is some
$I_{\bar x}$, and distinct $\bar x$ give distinct ideals (their quotients
are supported at distinct points of $X/\Gamma$).  This proves (2).
\end{proof}

\begin{lemma}\label{lem:varietypart}
Let $\Phi\in\Aut_{\mathrm{Lie}}(\M^\Gamma)$.  Then $\Phi$ permutes the
ideals $\{I_{\bar x}\}_{\bar x\in (X/\Gamma)(k)}$, inducing a bijection
$\bar\varphi\colon(X/\Gamma)(k)\to(X/\Gamma)(k)$ which is a variety
automorphism, and which coincides with the automorphism $\bar\varphi$ of
Lemma~\ref{lem:centroid}.  It lifts to a $\Gamma$-equivariant variety
automorphism $\varphi$ of $X$ normalizing the $\Gamma$-action.
\end{lemma}
\begin{proof}
$\Phi$ carries maximal ideals of codimension $\dim\g$ with simple quotient
to ideals of the same type, hence permutes $\{I_{\bar x}\}$ by
Lemma~\ref{lem:fibres}(2); define $\bar\varphi$ by
$\Phi(I_{\bar x})=I_{\bar\varphi(\bar x)}$.  Under the centroid
identification $C(\M^\Gamma)\cong A^\Gamma=\OO(X/\Gamma)$ of
Lemma~\ref{lem:centroid}, $I_{\bar x}=\mathfrak m_{\bar x}\M^\Gamma$ is the
ideal generated by the centroidal action of $\mathfrak m_{\bar x}$, so the
permutation $\bar\varphi$ agrees with the algebra automorphism
$\varphi^*=\Phi(\cdot)\Phi^{-1}$ of $\OO(X/\Gamma)$; thus $\bar\varphi$ is
a morphism of varieties with inverse $\overline{\Phi^{-1}}$, i.e.\ an
automorphism of $X/\Gamma$.  Finally, we must lift $\bar\varphi\in\Aut(X/\Gamma)$ to a
$\Gamma$-normalizing $\varphi\in\Aut(X)$.  This is \emph{not} automatic for
a general Galois cover: $\bar\varphi$ lifts iff $\bar\varphi^*$ pulls the
cover $X\to X/\Gamma$ back to an isomorphic $\Gamma$-cover, and there is in
general an obstruction (a class in $H^1(X/\Gamma,\underline{\Aut}\,\Gamma$-cover$)$).
This is the content of the standing hypothesis (H4) of \S\ref{sec:setup}.
Under (H4) such a lift $\varphi$ exists, with
$c(\gamma):=\varphi\gamma\varphi^{-1}\in\Gamma$ defining an automorphism
$c\in\Aut(\Gamma)$; the lift is unique up to a deck transformation in
$\Gamma$, an ambiguity absorbed into $\Psi$ below.  Hypothesis (H4) holds
in all the examples treated here: for the loop case $X=\Gm$, $X/\Gamma=\Gm$
and the cover is the $m$-th power map $t\mapsto t^m$, and every
$\bar\varphi\in\Aut(\Gm)=k^\times\rtimes\Z/2$ lifts to a monomial
automorphism of $\Gm$ (computed in \S\ref{sec:loop}); for the multiloop
case $X=(\Gm)^n$ the cover is $t_j\mapsto t_j^{m_j}$ and every admissible
monomial $\bar\varphi$ lifts to a monomial automorphism normalizing
$\Gamma$ (\S\ref{sec:multiloop}).  More generally (H4) holds whenever
$\mathrm{Pic}(X/\Gamma)$ and $H^1_{\mathrm{\acute et}}(X/\Gamma,\Gamma)$
vanish (e.g.\ $X/\Gamma$ a torus), by Pianzola's $H^1$-vanishing results
for Laurent-polynomial rings.
\end{proof}

\section{Reconstruction of the fibre data and proof of
Theorem~\ref{thm:main}}\label{sec:autthm}

\begin{lemma}[Fibrewise rigidity]\label{lem:fibrerigid}
Let $\Phi\in\Aut_{\mathrm{Lie}}(\M^\Gamma)$ with associated variety
automorphism $\varphi$ (Lemma~\ref{lem:varietypart}).  Then for each
$x\in X(k)$ there is a unique $\Psi(x)\in\Aut_{\mathrm{Lie}}(\g)$ such that
\begin{equation}\label{eq:fiber}
(\Phi f)(x)=\Psi(x)\bigl(f(\varphi^{-1}x)\bigr),\qquad f\in\M^\Gamma .
\end{equation}
The map $\Psi\colon X(k)\to\Aut_{\mathrm{Lie}}(\g)=G$ is regular (a
morphism $X\to G$), and it satisfies the equivariance condition
\textnormal{\eqref{eq:equiv}} proved below, namely
$\Psi(c(\gamma)x)\,\gamma=c(\gamma)\,\Psi(x)$ as elements of $G$, where
$c(\gamma)=\varphi\gamma\varphi^{-1}\in\Gamma$ and $\gamma,\,c(\gamma)$
denote the images in $G=\Aut(\g)$ of the corresponding elements of
$\Gamma$.  In particular $\Psi$ lies in the twisted gauge group
$\mathcal G_\Gamma(X,\g)$ of \eqref{eq:gaugegroup} (a regular $G$-valued
map, not necessarily $N_\Gamma(\g)$-valued), and $(\varphi,\Psi)$ is a
$\Gamma$-equivariant twisted automorphism in the sense of
Definition~\ref{def:emap}.
\end{lemma}
\begin{proof}
Fix $x\in X(k)$ and let $\bar x$ be its image.  Since
$\Phi(I_{\varphi^{-1}\bar x})=I_{\bar x}$ (Lemma~\ref{lem:varietypart},
applied to $\Phi^{-1}$ and rearranged), $\Phi$ descends to a Lie algebra
isomorphism of the simple quotients
\[
\bar\Phi_x\colon \M^\Gamma/I_{\varphi^{-1}\bar x}\;\xrightarrow{\ \sim\ }\;
\M^\Gamma/I_{\bar x},
\]
which under the identifications of Lemma~\ref{lem:fibres}(1)
($\mathrm{ev}_{\varphi^{-1}x}$ on the source, $\mathrm{ev}_x$ on the
target) becomes an automorphism $\Psi(x)$ of the simple Lie algebra $\g$.
This is \eqref{eq:fiber}, and $\Psi(x)$ is unique because
$\mathrm{ev}_x|_{\M^\Gamma}$ is surjective onto $\g$
(Lemma~\ref{lem:fibres}(1)).

\emph{Regularity.}  Fix a $k$-basis $u_1,\dots,u_d$ of $\g$.  By
Lemma~\ref{lem:fibres}(1), for each $j$ and each $\bar x$ there is
$f_j\in\M^\Gamma$ with $f_j(\varphi^{-1}x)=u_j$; the matrix of $\Psi(x)$ in
the basis $(u_j)$ is $\bigl[(\Phi f_j)(x)\bigr]_j$ expressed in $(u_i)$.
The entries $x\mapsto(\Phi f_j)(x)$ are components of the regular map
$\Phi f_j\in\M=\OO(X)\otimes\g$, hence regular functions of $x$ on the
locus where $f_j(\varphi^{-1}\cdot)$ stays a basis; covering $X$ by such
loci (possible because $\mathrm{ev}$ is surjective on $\M^\Gamma$, an open
condition) shows $\Psi\colon X\to\GL(\g)$ is a morphism, and its image lies
in the closed subgroup $G=\Aut(\g)\subset\GL(\g)$, so $\Psi\colon X\to G$
is a morphism of varieties.

\emph{Equivariance.}  Write $c\in\Aut(\Gamma)$ for the automorphism with
$\varphi\,\gamma\,\varphi^{-1}=c(\gamma)$ (Lemma~\ref{lem:varietypart}).
Conjugating this relation gives the two equivalent forms we shall use:
\begin{equation}\label{eq:phicomm}
\varphi^{-1}\,c(\gamma)=\gamma\,\varphi^{-1},
\qquad\text{equivalently}\qquad
\varphi^{-1}\,\delta=c^{-1}(\delta)\,\varphi^{-1}\ \ (\delta\in\Gamma).
\end{equation}
Now fix $x$ and $\gamma$, and put $\delta:=c(\gamma)\in\Gamma$.  Evaluate
\eqref{eq:fiber} at the point $\delta x$:
\begin{equation}\label{eq:e1}
(\Phi f)(\delta x)=\Psi(\delta x)\,\bigl(f(\varphi^{-1}\delta x)\bigr)
=\Psi(\delta x)\,\bigl(f(\gamma\,\varphi^{-1}x)\bigr),
\end{equation}
where the second equality uses $\varphi^{-1}\delta=\gamma\varphi^{-1}$ from
\eqref{eq:phicomm}.  Since $f\in\M^\Gamma$, i.e.\ $f(\gamma y)=
\gamma\bigl(f(y)\bigr)$ for all $y$, the right side of \eqref{eq:e1}
equals
\begin{equation}\label{eq:e2}
\Psi(\delta x)\,\gamma\,\bigl(f(\varphi^{-1}x)\bigr),
\end{equation}
where $\gamma$ now denotes the action of $\gamma\in\Gamma$ on $\g$, i.e.\
its image in $G=\Aut(\g)$.  On the other hand $\Phi f\in\M^\Gamma$, so it
too is $\Gamma$-invariant, and since $\delta\in\Gamma$,
\begin{equation}\label{eq:e3}
(\Phi f)(\delta x)=\delta\,\bigl((\Phi f)(x)\bigr)
=\delta\,\Psi(x)\,\bigl(f(\varphi^{-1}x)\bigr),
\end{equation}
where $\delta$ denotes the action of $\delta\in\Gamma$ on $\g$.  As $f$
ranges over $\M^\Gamma$, the vector $f(\varphi^{-1}x)$ ranges over all of
$\g$ (Lemma~\ref{lem:fibres}(1)); equating \eqref{eq:e2} and \eqref{eq:e3}
therefore yields the identity of endomorphisms of $\g$
\[
\Psi(\delta x)\,\gamma=\delta\,\Psi(x),
\qquad\text{i.e.}\qquad
\Psi\bigl(c(\gamma)\,x\bigr)\,\gamma=c(\gamma)\,\Psi(x),
\]
where on the left $\gamma\in G$ and on the right $c(\gamma)\in G$ via
$\Gamma_\g$.  Equivalently,
\begin{equation}\label{eq:equiv}
\Psi\bigl(c(\gamma)\,x\bigr)=c(\gamma)\,\Psi(x)\,\gamma^{-1}\qquad
(\gamma\in\Gamma,\ x\in X(k)).
\end{equation}
This is exactly the twisted equivariance \eqref{eq:equivdef} of
Definition~\ref{def:emap}.  We emphasize (correcting an earlier draft) that
\eqref{eq:equiv} does \emph{not} force $\Psi(x)\in N_\Gamma(\g)$: it only
relates $\Psi$ at the distinct points $x$ and $c(\gamma)x$ of an orbit.
Indeed, solving $\Psi(x)=c(\gamma)^{-1}\Psi(c(\gamma)x)\gamma$ on a free
orbit $\{x_1,\dots,x_{|\Gamma|}\}$ (say with $c=\id$) gives
$\Psi(\gamma x_1)=\gamma\,\Psi(x_1)\,\gamma^{-1}$, leaving $\Psi(x_1)\in G$
\emph{free}.  Thus $\Psi$ is a regular section of the (constant) group
scheme $G$ over $X$ satisfying \eqref{eq:equiv}, i.e.\ an element of the
twisted gauge group $\mathcal G_\Gamma(X,\g)$ of \eqref{eq:gaugegroup}, and
$(\varphi,\Psi)$ is a $\Gamma$-equivariant twisted automorphism in the
sense of Definition~\ref{def:emap}.
\end{proof}

\begin{proof}[Proof of Theorem~\ref{thm:main}]
\emph{Induced operators are automorphisms.}  Let $(\varphi,\Psi)$ be a
$\Gamma$-equivariant twisted automorphism and $\Phi$ the operator
\eqref{eq:Phi}.  For $f,g\in\M$ and $x\in X(k)$,
\[
(\Phi[f,g])(x)=\Psi(x)\bigl([f,g](\varphi^{-1}x)\bigr)
=\Psi(x)\bigl[f(\varphi^{-1}x),g(\varphi^{-1}x)\bigr]
=\bigl[\Psi(x)f(\varphi^{-1}x),\Psi(x)g(\varphi^{-1}x)\bigr]
=[\Phi f,\Phi g](x),
\]
where the third equality is because $\Psi(x)\in\Aut_{\mathrm{Lie}}(\g)$.
Thus $\Phi$ is a Lie endomorphism of $\M$; it is bijective with inverse
induced by $(\varphi^{-1},x\mapsto\Psi(\varphi^{-1}x)^{-1})$.  By
hypothesis (Definition~\ref{def:emap}) $\Phi$ preserves $\M^\Gamma$, so
$\Phi|_{\M^\Gamma}\in\Aut_{\mathrm{Lie}}(\M^\Gamma)$.  Therefore
$\mathcal E(X,\g)^\Gamma\subseteq\Aut(\M^\Gamma)$.

\emph{Every automorphism is induced.}  Let
$\Phi\in\Aut(\M^\Gamma)$.  Lemma~\ref{lem:varietypart} produces the variety
part $\varphi$ (normalizing $\Gamma$), and Lemma~\ref{lem:fibrerigid}
produces the regular map $\Psi\colon X\to G=\Aut(\g)$ satisfying
the twisted equivariance \eqref{eq:equiv} together with the pointwise
identity \eqref{eq:fiber} (recall that $\Psi$ is
valued in the \emph{full} automorphism group $G$, not merely in
$N_\Gamma(\g)$).
Identity \eqref{eq:fiber} says exactly that $\Phi$ agrees on $\M^\Gamma$
with the induced operator of \eqref{eq:Phi}.  By construction
$(\varphi,\Psi)$ preserves $\M^\Gamma$ (it equals $\Phi$ there), so it is a
$\Gamma$-equivariant twisted automorphism and
$\Phi\in\mathcal E(X,\g)^\Gamma$.  Uniqueness of $(\varphi,\Psi)$: $\varphi$
is determined by the action on $C(\M^\Gamma)=A^\Gamma$
(Lemma~\ref{lem:centroid}), and $\Psi$ is then determined pointwise by
\eqref{eq:fiber} and the surjectivity of $\mathrm{ev}_x|_{\M^\Gamma}$
(Lemma~\ref{lem:fibres}).  Hence
$\Aut(\M^\Gamma)=\mathcal E(X,\g)^\Gamma$.
\end{proof}

\begin{remark}[Group structure of $\mathcal E(X,\g)^\Gamma$]
\label{rem:groupstr}
Composition of induced operators gives
$(\varphi_2,\Psi_2)\circ(\varphi_1,\Psi_1)
=(\varphi_2\varphi_1,\ x\mapsto\Psi_2(x)\,\Psi_1(\varphi_2^{-1}x))$,
which is the multiplication in the semidirect product
\[
\mathcal E(X,\g)^\Gamma\;\cong\;
\Aut(X,\Gamma)\ \ltimes\ \mathcal G_\Gamma(X,\g),
\]
where $\Aut(X,\Gamma)$ is the group of variety automorphisms of $X$
normalizing the $\Gamma$-action, and $\mathcal G_\Gamma(X,\g)$ is the
twisted equivariant gauge group \eqref{eq:gaugegroup} of regular maps
$\Psi\colon X\to G=\Aut(\g)$ satisfying \eqref{eq:equiv}, with
$\Aut(X,\Gamma)$ acting by $\varphi\cdot\Psi=\Psi\circ\varphi^{-1}$ and
$c\in\Aut(\Gamma)$ twisting the equivariance.  This is precisely the
content of Theorem~\ref{thm:main} in group-theoretic form.

\medskip
\noindent\emph{Structure of $\mathcal G_\Gamma(X,\g)$.}  Fix a set
$R\subseteq X(k)$ of representatives of the $\Gamma$-orbits (i.e.\ a section
of $X(k)\to (X/\Gamma)(k)$).  By \eqref{eq:equiv} a gauge map $\Psi$ is
determined by its restriction $\Psi|_R$, which may be \emph{any} regular
$G$-valued function on $R$ (no constraint), and conversely any such
restriction extends uniquely by $\Psi(c(\gamma)x)=c(\gamma)\Psi(x)
\gamma^{-1}$.  Thus, ignoring regularity across orbits, $\mathcal
G_\Gamma(X,\g)$ is a ``twisted form'' of the full gauge group
$\mathrm{Mor}(X/\Gamma,G)$; equivalently it is the group of regular
sections of the group scheme $\underline G=(\pi_*G_X)^\Gamma$ over
$X/\Gamma$ obtained by Galois descent of the constant group $G$ along the
$\Gamma$-cover $\pi\colon X\to X/\Gamma$ twisted by $c$.  Its component
group (the constant/$\pi_0$ part) is governed by $N_\Gamma(\g)$, which is
where the diagram/outer symmetries live; this is what enters the explicit
loop and multiloop formulas below.
\end{remark}

\section{Twisted loop algebras: proof of Theorem~\ref{thm:loop}}
\label{sec:loop}

Here $X=\Gm=\Spec k[t,t^{-1}]$, $A=k[t,t^{-1}]$, $\g$ simple,
$\Gamma=\Z/m=\langle\gamma\rangle$ acting by $\sigma=\gamma\cdot$ of order
$m$ on $\g$ and by $\gamma\cdot t=\zeta t$ (so $\gamma^{-1}\cdot t=
\zeta^{-1}t$ on functions: $(\gamma\cdot p)(t)=p(\zeta^{-1}t)$).  We verify
(H1)--(H3): $k$ algebraically closed, $A$ finitely generated and reduced
(H1); $\Gamma$ finite and the action $t\mapsto\zeta^{j}t$ on
$\Gm(k)=k^\times$ is free because $\zeta^j t=t$ with $t\ne0$ forces
$\zeta^j=1$, i.e.\ $j\equiv0$ (H2); $\g$ simple (H3).  Theorem~\ref{thm:main}
applies.

\subsection{Identification with the standard loop algebra}
Decompose $\g=\bigoplus_{j\in\Z/m}\g_{\bar j}$ into $\sigma$-eigenspaces,
$\sigma|_{\g_{\bar j}}=\zeta^{j}$.  A Laurent series $f=\sum_i a_i t^i$
($a_i\in\g$) is $\Gamma$-invariant iff
$\sigma(a_i)\zeta^{-i}=a_i$ for all $i$ (computed in
\S\ref{sec:setup}), i.e.\ $a_i\in\g_{\bar i}$.  Hence
\[
\M^\Gamma=L(\g,\sigma)=\bigoplus_{i\in\Z}\g_{\bar i}\,t^{i},
\]
the standard twisted loop algebra.

\subsection{The variety automorphisms $\Aut(\Gm,\Gamma)$}
$\Aut_{\mathrm{var}}(\Gm)=\{t\mapsto at^{\pm1}:a\in k^\times\}\cong
k^\times\rtimes\Z/2$.  Such $\varphi$ normalizes the $\Gamma$-action
$t\mapsto\zeta t$ iff $\varphi\circ(t\mapsto\zeta t)\circ\varphi^{-1}\in
\Gamma$.  For $\varphi\colon t\mapsto at$: $\varphi(\zeta\varphi^{-1}(t))
=\varphi(\zeta a^{-1}t)=\zeta t$, which is $\gamma$ again, so $c=\id$.  For
$\varphi\colon t\mapsto at^{-1}$:
$\varphi(\zeta\varphi^{-1}(t))=\varphi(\zeta a t^{-1})
=a(\zeta a t^{-1})^{-1}=\zeta^{-1}t$, which is $\gamma^{-1}$, so
$c(\gamma)=\gamma^{-1}$ (inversion of $\Gamma$).  Thus every variety
automorphism of $\Gm$ normalizes $\Gamma$, and
\[
\Aut(\Gm,\Gamma)=k^\times\rtimes\Z/2,
\]
where the $\Z/2$ (inversion $\iota\colon t\mapsto t^{-1}$) induces
$c(\gamma)=\gamma^{-1}$, hence on $\g$ forces conjugation taking
$\sigma\mapsto\sigma^{-1}$.

\subsection{The gauge maps and the constraint \eqref{eq:equiv}}
By Theorem~\ref{thm:main} and Remark~\ref{rem:groupstr},
\[
\Aut(L(\g,\sigma))\;\cong\;
(k^\times\rtimes\Z/2)\ \ltimes\ \mathcal G_\Gamma(\Gm,\g).
\]
We compute the twisted equivariant gauge group $\mathcal G_\Gamma(\Gm,\g)$
for $c=\id$ (the orientation $t\mapsto at$ piece; the inversion piece is
handled by the $\Z/2$).  A morphism $\Psi\colon\Gm\to G=\Aut(\g)$ is a
$G$-valued Laurent loop; with $c=\id$ and the $\Gamma$-action
$\gamma\cdot x=\zeta^{-1}x$, condition \eqref{eq:equiv} reads
\[
\Psi(\zeta^{-1} x)=\gamma\,\Psi(x)\,\gamma^{-1}=\sigma\,\Psi(x)\,\sigma^{-1},
\qquad x\in k^\times,
\]
equivalently $\Psi(\zeta x)=\sigma^{-1}\Psi(x)\sigma$; i.e.\ $\Psi$ is
equivariant for the order-$m$ rotation of $\Gm$ and the conjugation action
of $\langle\sigma\rangle$ on $G$.

\emph{Constant loops: centralizer vs.\ normalizer bookkeeping.}  A
\emph{constant} gauge $\Psi\equiv\theta$ paired with a variety part
$\varphi$ inducing $c\in\Aut(\Gamma)$ is admissible (satisfies
\eqref{eq:equiv}) iff for the generator $\gamma$ with image $\sigma\in G$,
\begin{equation}\label{eq:loopconstconstraint}
\theta=c(\gamma)\,\theta\,\gamma^{-1}
\quad\Longleftrightarrow\quad
\theta\,\sigma\,\theta^{-1}=\sigma_{c(\gamma)},
\end{equation}
where $\sigma_{c(\gamma)}$ is the image of $c(\gamma)$ in $G$.  For
$\Gamma=\Z/m$ the only $c\in\Aut(\Gamma)$ realized by a variety
automorphism of $\Gm$ (\S\ref{sec:loop}, ``variety automorphisms'') are
$c=\id$ (from rescalings $t\mapsto at$, including $a=1$) and
$c\colon\gamma\mapsto\gamma^{-1}$ (from the inversion $\iota$).  Hence:
\begin{itemize}
\item with the rescaling part ($c=\id$): \eqref{eq:loopconstconstraint}
reads $\theta\sigma\theta^{-1}=\sigma$, i.e.\ $\theta\in Z_G(\sigma)$, the
\emph{centralizer}; these constant gauges pair with any rescaling
$\varphi_a$;
\item with the inversion ($c\colon\gamma\mapsto\gamma^{-1}$):
\eqref{eq:loopconstconstraint} reads $\theta\sigma\theta^{-1}=\sigma^{-1}$;
such $\theta$ exist iff $\sigma$ and $\sigma^{-1}$ are conjugate in
$\Aut(\g)$ (always true for diagram automorphisms, as $\sigma$ and
$\sigma^{-1}$ are conjugate diagram automorphisms), and they pair with the
inversion $\iota$ (consistent with $L(\g,\sigma)\cong L(\g,\sigma^{-1})$).
\end{itemize}
Thus the set of \emph{constant} gauges admissible together with \emph{some}
variety part is exactly
\[
\{\theta\in G:\theta\sigma\theta^{-1}\in\{\sigma,\sigma^{-1}\}\}
=Z_G(\sigma)\,\cup\,\{\theta:\theta\sigma\theta^{-1}=\sigma^{-1}\},
\]
which is the subgroup $N^{\pm}_\Gamma(\g):=\{\theta:\theta\langle\sigma
\rangle\theta^{-1}=\langle\sigma\rangle,\ \theta\sigma\theta^{-1}\in\{\sigma,
\sigma^{-1}\}\}\subseteq N_\Gamma(\g)$.  (For $m\le2$, $N^\pm_\Gamma(\g)=
N_\Gamma(\g)$ since then every power of $\sigma$ in $\langle\sigma\rangle$ is
$\sigma^{\pm1}$.  For $m\ge3$, an element $\theta\in N_\Gamma(\g)$ with
$\theta\sigma\theta^{-1}=\sigma^{r}$, $r\not\equiv\pm1$, is \emph{not}
realizable by a \emph{constant} gauge on $\Gm$, precisely because $\Gm$
admits no variety automorphism inducing $\gamma\mapsto\gamma^{r}$ for such
$r$; this is the honest statement of item (1) of Theorem~\ref{thm:loop} and
corrects the earlier over-claim that all of $N_\Gamma(\g)$ arises from
constant gauges.)  The remaining elements of $N_\Gamma(\g)$, when they
exist, are realized only by \emph{non-constant} gauges or not at all, and do
not enlarge $\Aut(L)$ beyond the group computed here.

\emph{Decomposition of the gauge group into inner and outer parts.}
Decompose $G=\Aut(\g)$ with identity component $G^\circ=\mathrm{Inn}(\g)$
(a connected adjoint semisimple group, since $\g$ is simple) and finite
component group $\mathrm{Out}(\g)=G/G^\circ$.  Since $\Gm$ is connected,
its image under any morphism $\Psi\colon\Gm\to G$ lies in a single
connected component $\theta_0 G^\circ$ of $G$; choosing any constant
representative $\theta_0\in N_\Gamma(\g)$ of that component we may write
$\Psi=\theta_0\cdot\Psi^\circ$ with $\Psi^\circ:=\theta_0^{-1}\Psi\colon
\Gm\to G^\circ=\mathrm{Inn}(\g)$ a loop into the inner group.  Thus
\begin{equation}\label{eq:gaugesplit}
\mathcal G_\Gamma(\Gm,\g)=
\mathcal G_\Gamma(\Gm,\g)^\circ\cdot R,
\end{equation}
where $\mathcal G_\Gamma(\Gm,\g)^\circ:=\{\Psi\in\mathcal G_\Gamma(\Gm,\g):
\Psi(\Gm)\subseteq G^\circ\}$ is the subgroup of inner gauge loops and $R$
is a finite set of constant representatives $\theta_0$ of the components
$\theta_0 G^\circ$ that admit an equivariant constant gauge, i.e.\ with
$\theta_0\in N_\Gamma(\g)$ (so that $\theta_0\equiv$const satisfies
\eqref{eq:equiv} after pairing with a suitable $\varphi$).  The second
factor is the finite outer/diagram part (item (1)), the first the
connected (inner) gauge loops.

\emph{The inner gauge loops are exactly the inner automorphisms of $L$.}
We make this identification precise, distinguishing the two groups the
critic rightly warns must not be conflated:
\begin{itemize}
\item $\mathcal G^\circ:=\{\Psi\in\mathcal G_\Gamma(\Gm,\g):
\Psi(\Gm)\subseteq G^\circ=\mathrm{Inn}(\g)\}$, the $\varphi=\id$,
$\mathrm{Inn}(\g)$-valued, $\sigma$-equivariant gauge loops (a subgroup of
$\Aut(L)$ by Theorem~\ref{thm:main});
\item $\mathrm{Inn}(L)$, the subgroup of $\Aut(L)$ generated by the
automorphisms $\exp(\mathrm{ad}\,\xi)$ with $\xi\in L$ \emph{locally
ad-nilpotent} (these exponentials are well defined because $\mathrm{ad}\,
\xi$ acts locally nilpotently on $L$ for such $\xi$).
\end{itemize}
We prove $\mathcal G^\circ=\mathrm{Inn}(L)$.

\emph{(i) $\mathrm{Inn}(L)\subseteq\mathcal G^\circ$.}  It suffices to treat
the generators.  Let $v\in\g_{\bar i}$ be ad-nilpotent in $\g$ and
$\xi=v\,t^{i}\in L$.  Then $(\mathrm{ad}\,\xi)^N f$ raises the degree by
$Ni$ while applying $(\mathrm{ad}\,v)^N$ to coefficients, so for fixed
$f\in L$ (finitely many nonzero homogeneous components) only finitely many
terms survive; hence $\mathrm{ad}\,\xi$ is locally nilpotent and
$\exp(\mathrm{ad}\,\xi)$ is defined.  Its action is
$(\exp(\mathrm{ad}\,\xi)f)(t)=\exp(t^i\,\mathrm{ad}\,v)\,f(t)=\Psi_\xi(t)
f(t)$, with $\Psi_\xi(t)=\exp(t^i\,\mathrm{ad}\,v)\in\mathrm{Inn}(\g)$, an
$\mathrm{Inn}(\g)$-valued gauge loop with $\varphi=\id$.  It is
$\sigma$-equivariant: since $\sigma v=\zeta^{i}v$,
\[
\Psi_\xi(\zeta x)=\exp(\zeta^i x^i\,\mathrm{ad}\,v)
=\exp(x^i\,\mathrm{ad}(\sigma v))
=\sigma\exp(x^i\,\mathrm{ad}\,v)\sigma^{-1}=\sigma\,\Psi_\xi(x)\,\sigma^{-1},
\]
which is \eqref{eq:equiv} for $c=\id$.  Thus $\Psi_\xi\in\mathcal G^\circ$,
and products of such generate a subgroup of $\mathcal G^\circ$; so
$\mathrm{Inn}(L)\subseteq\mathcal G^\circ$.

\emph{(ii) $\mathcal G^\circ\subseteq\mathrm{Inn}(L)$.}  This is the
deepest point of the loop computation, and we are careful to separate what we
prove by hand from the one external generation theorem we invoke.  We first
record the structure we need over the untwisted Laurent ring.

\begin{lemma}[Elementary generation contains the torus]\label{lem:steinberg}
Let $\bar G=\mathrm{Inn}(\g)$ be the adjoint Chevalley group of the simple
$\g$, and for a commutative $k$-algebra $R$ let
\[
E(R):=\bigl\langle\, x_\alpha(p)=\exp\bigl(p\,\mathrm{ad}\,e_\alpha\bigr)\ :\
\alpha\in\Delta,\ p\in R\,\bigr\rangle\subseteq\bar G(R)
\]
be the elementary subgroup generated by the root subgroups (here $\Delta$ is
the root system relative to a fixed Cartan $\mathfrak h\subseteq\g$ and
$e_\alpha\in\g_\alpha$ a fixed root vector).  Then for every unit $u\in
R^\times$ and every $\alpha\in\Delta$ the semisimple torus elements
\[
w_\alpha(u):=x_\alpha(u)\,x_{-\alpha}(-u^{-1})\,x_\alpha(u),\qquad
h_\alpha(u):=w_\alpha(u)\,w_\alpha(-1)
\]
lie in $E(R)$; consequently $E(R)$ contains the full split maximal torus
$\bar T(R)=\langle h_\alpha(u):\alpha\in\Delta,\ u\in R^\times\rangle$ of
$\bar G(R)$.
\end{lemma}
\begin{proof}
The displayed identities are the standard Steinberg relations valid in any
Chevalley group over any commutative ring (Steinberg, \emph{Lectures on
Chevalley groups}, \S3; the formulas for $w_\alpha(u),h_\alpha(u)$ make sense
because $u$ is a unit).  Each factor $x_{\pm\alpha}(\cdot)$ lies in $E(R)$ by
definition, hence so do $w_\alpha(u)$ and $h_\alpha(u)$.  For the adjoint
group the elements $h_\alpha(u)$ generate the split torus $\bar T(R)$ (the
cocharacter lattice of $\bar T$ is the coroot lattice, which is spanned by
the $\alpha^\vee$), proving the last claim.
\end{proof}

We will use $R=A=k[t^{\pm1}]$.  The crucial input is the following
\emph{generation theorem}, which we cite (it is not elementary and is the one
external ingredient used in the proof of the equality $\mathcal G^\circ=\mathrm{Inn}(L)$).

\begin{theorem}[Generation of loop groups; {Pianzola; Gille--Pianzola; Peterson--Kac}]
\label{thm:looptorgen}
Let $\bar G=\mathrm{Inn}(\g)$ be adjoint simple, $A=k[t^{\pm1}]$, and let
$\sigma\in\bar G(k)$ (more generally $\sigma\in\Aut(\g)$) be of finite order
$m$ acting on $\bar G(A)$ by the twisted rotation
$(\gamma * g)(t)=\sigma\,g(\zeta^{-1}t)\,\sigma^{-1}$, with fixed-point group
$\mathcal G^\circ=\bar G(A)^{\langle\gamma*\rangle}$.  Then $\mathcal
G^\circ$ is generated by its equivariant root subgroups, i.e.\ by the
elements $\exp(\mathrm{ad}\,\xi)$ with $\xi\in L=L(\g,\sigma)$ a homogeneous
\textnormal{ad}-nilpotent element (so $\xi=v\,t^i$ with $v\in\g_{\bar i}$ and
$v$ \textnormal{ad}-nilpotent in $\g$), together with the equivariant torus
elements arising from \textnormal{Lemma~\ref{lem:steinberg}}.
\end{theorem}

\noindent\emph{On the status of Theorem~\ref{thm:looptorgen}.}  This is a
genuine theorem and we do \emph{not} reprove it.  Its content is twofold.
(a) Over $A=k[t^{\pm1}]$ the untwisted adjoint loop group $\bar G(A)$ is
generated by its root subgroups together with $\bar T(A)$: for $\g$ of rank
$\ge2$ this is unstable $K_1$-triviality for split groups over the regular
$1$-dimensional ring $A$ (Suslin--Kopeiko, Abe), while for $\g=\mathfrak{sl}_2$
the apparent failure ``$E_2\ne SL_2$ over $k[t^{\pm1}]$'' is repaired exactly
by adjoining the torus $\bar T(A)$ supplied by Lemma~\ref{lem:steinberg} and
by passing to the adjoint group, so that $\bar G(A)=E(A)\cdot\bar T(A)$ holds
in all ranks.  (b) The fixed-point (twisted) group $\mathcal G^\circ$ is the
group of $A^\Gamma$-points of the adjoint twisted form $\underline{\bar G}$
over $A^\Gamma=k[s^{\pm1}]$ ($s=t^m$) obtained by Galois descent along the
$\Gamma$-cover; by Pianzola's vanishing $H^1(A^\Gamma,\underline{\bar G})=0$
for loop-reductive groups over Laurent rings, $\underline{\bar G}(A^\Gamma)$
is generated by its relative root subgroups, and each relative root subgroup
element is precisely an equivariant exponential $\exp(\mathrm{ad}\,\xi)$ of
the stated form.  We give precise references in the bibliography.

We now finish (ii) by hand, given Theorem~\ref{thm:looptorgen}.  Let
$\Psi\in\mathcal G^\circ$.  By Theorem~\ref{thm:looptorgen}, $\Psi$ is a
finite product of two kinds of generators:
\begin{enumerate}
\item[(a)] equivariant exponentials $\exp(\mathrm{ad}\,\xi)$,
$\xi=v\,t^i\in L$ with $v\in\g_{\bar i}$ ad-nilpotent; these lie in
$\mathrm{Inn}(L)$ by definition;
\item[(b)] equivariant torus elements coming from Lemma~\ref{lem:steinberg}.
\end{enumerate}
For (b) we must check that the relevant torus elements again lie in
$\mathrm{Inn}(L)=\langle\exp(\mathrm{ad}\,\xi)\rangle$.  This is immediate
from Lemma~\ref{lem:steinberg}: each $h_\alpha(u)$ and each $w_\alpha(u)$ is,
by the Steinberg formulas, a finite product of root-subgroup elements
$x_{\pm\alpha}(p)=\exp(p\,\mathrm{ad}\,e_{\pm\alpha})$ with $p\,e_{\pm\alpha}$
ad-nilpotent; and in the equivariant setting one uses the
$\Gamma$-symmetrised generators
\[
\widetilde x_\alpha(p):=\prod_{j=0}^{m-1}\gamma^{j}\!*\!x_\alpha(p),
\qquad
\gamma\!*\!x_\alpha(p)
=\sigma\,\exp\!\bigl(p(\zeta^{-1}t)\,\mathrm{ad}\,e_\alpha\bigr)\,\sigma^{-1}
=x_{\sigma\alpha}\bigl(c_\alpha\,p(\zeta^{-1}t)\bigr),
\]
where $\sigma e_\alpha=c_\alpha e_{\sigma\alpha}$.  Each factor
$\gamma^{j}\!*\!x_\alpha(p)=\exp(\mathrm{ad}\,\eta_j)$ with $\eta_j=
p(\zeta^{-j}t)\,\sigma^{j}e_\alpha$ a homogeneous ad-nilpotent element of the
\emph{untwisted} loop algebra $\M=\bigoplus_i\g\,t^i$, hence lies in
$\mathrm{Inn}(\M)$, and the $\Gamma$-symmetrised product $\widetilde
x_\alpha(p)$ is by construction $\langle\gamma*\rangle$-invariant, so it lies
in $\mathcal G^\circ\subseteq\Aut(L)$ and is a finite product of
exponentials of ad-nilpotents \emph{that preserve $L$}; therefore
$\widetilde x_\alpha(p)\in\mathrm{Inn}(L)$.  The Steinberg combinations
$w_\alpha,h_\alpha$ built from these equivariant root elements thus also lie
in $\mathrm{Inn}(L)$.  Since every generator of $\mathcal G^\circ$ in
Theorem~\ref{thm:looptorgen} is of type (a) or is such a Steinberg
combination of type-(b) elements, we conclude $\Psi\in\mathrm{Inn}(L)$, i.e.\
$\mathcal G^\circ\subseteq\mathrm{Inn}(L)$.

\medskip
\noindent\emph{Honest scope of the loop-group input.}  We have reduced the equality
$\mathcal G^\circ=\mathrm{Inn}(L)$ to the single cited
Theorem~\ref{thm:looptorgen} (generation of the twisted loop group by its
equivariant root and torus subgroups) and then proved both inclusions by
hand from it.  We do not claim a new proof of
Theorem~\ref{thm:looptorgen}; it is a known result of loop-group theory
(Peterson--Kac generation by real root groups in the Kac--Moody setting;
Pianzola and Gille--Pianzola via $H^1$-vanishing and Galois descent for
loop-reductive groups over Laurent rings).  This is the correct and
literature-aligned resolution: the statement $\mathcal G^\circ=\mathrm{Inn}(L)$
is true but is \emph{not} elementary, and its depth resides exactly in
Theorem~\ref{thm:looptorgen}.

Combining (i) and (ii),
\begin{equation}\label{eq:innerloop}
\mathcal G^\circ=\mathrm{Mor}_\Gamma\bigl(\Gm,\mathrm{Inn}(\g)\bigr)
=\mathrm{Inn}(L).
\end{equation}
Combining \eqref{eq:gaugesplit} and \eqref{eq:innerloop}, the gauge group
$\mathcal G_\Gamma(\Gm,\g)$ is generated by $\mathrm{Inn}(L)$ together with
the finite group of constant outer symmetries (the image of $N_\Gamma(\g)$
in $\mathrm{Out}(\g)$); this is item (1) and item (4) of
Theorem~\ref{thm:loop}.

\subsection{Assembling Theorem~\ref{thm:loop}}
Collecting the variety part $k^\times\rtimes\Z/2$ and the gauge part, every
$\Phi\in\Aut(L)$ has the form \eqref{eq:loopaut}
$(\Phi f)(t)=\Psi(t)\,f(at^{\pm1})$.  The semidirect/extension structure is
exactly as displayed in Theorem~\ref{thm:loop}:
\begin{itemize}
\item the rescalings $\varphi_a:t\mapsto at$ ($a\in k^\times$) act by
$f(t)\mapsto f(a^{-1}t)$ and give the $k^\times$;
\item the inversion $\iota:t\mapsto t^{-1}$ gives the $\Z/2$, and forces
$\sigma\mapsto\sigma^{-1}$ on $\g$ (computed above), consistent with the
isomorphism $L(\g,\sigma)\cong L(\g,\sigma^{-1})$;
\item the constant outer/diagram part $\theta\in N^\pm_\Gamma(\g)$
(centralizer of $\sigma$ paired with rescalings, and the
$\sigma\mapsto\sigma^{-1}$ part paired with the inversion) realizes the
\emph{realizable} symmetries of $(\g,\langle\sigma\rangle)$, as analysed in
the constant-loop bookkeeping above;
\item the inner loops are $\mathrm{Inn}(L)=\mathcal G^\circ$ by
\eqref{eq:innerloop}.
\end{itemize}
The full automorphism group is, by Theorem~\ref{thm:main} and
Remark~\ref{rem:groupstr}, the semidirect product
\[
\Aut(L)\;\cong\;(k^\times\rtimes\Z/2)\ \ltimes\ \mathcal G_\Gamma(\Gm,\g),
\qquad
\mathcal G_\Gamma(\Gm,\g)=\mathcal G^\circ\cdot R
=\mathrm{Inn}(L)\cdot R,
\]
where $\mathcal G_\Gamma(\Gm,\g)$ is the \emph{$G$-valued} twisted gauge
group \eqref{eq:gaugegroup} (\emph{not} $N_\Gamma(\g)$-valued) and $R$ is the finite set of component representatives
$\theta_0\in N_\Gamma(\g)$ from \eqref{eq:gaugesplit}.  Passing to outer
automorphisms (quotient by $\mathrm{Inn}(L)$ and by the variety group
$k^\times\rtimes\Z/2$) leaves exactly the finite group of diagram symmetries
of the pair $(\g,\langle\sigma\rangle)$, the image of $N_\Gamma(\g)$ in
$\mathrm{Out}(\g)$, namely
\[
\mathrm{Out}(L):=\Aut(L)/\bigl((k^\times\rtimes\Z/2)\cdot\mathrm{Inn}(L)\bigr)
\;\cong\;N_\Gamma(\g)\big/\bigl(N_\Gamma(\g)\cap
(\mathrm{Inn}(\g)\!\cdot\!\langle\sigma\rangle)\bigr),
\]
which maps to the stabiliser
$\mathrm{Stab}(\bar\sigma)\subseteq\mathrm{Out}(\g)$ of the conjugacy class
$\bar\sigma$ of $\sigma$ in $\mathrm{Out}(\g)$.  This matches the known
description of the automorphism group of the affine Kac--Moody algebra
$X_N^{(r)}$ (Kac, \emph{Infinite-dimensional Lie algebras}, Ch.~8;
Peterson--Kac), with the centre/derivation extension stripped off, since
$L=[\hat{\mathfrak g},\hat{\mathfrak g}]/\text{centre}$ is the loop
quotient.  This is the explicit description claimed in
Theorem~\ref{thm:loop}, the loop specialization of
Remark~\ref{rem:groupstr}.

\section{Multiloop algebras: proof of Theorem~\ref{thm:multiloop}}
\label{sec:multiloop}

Here $X=(\Gm)^n=\Spec k[t_1^{\pm1},\dots,t_n^{\pm1}]$,
$\Gamma=\prod_{j=1}^n\Z/m_j=\langle\gamma_1,\dots,\gamma_n\rangle$, with
$\gamma_j\cdot t_i=\zeta_{m_j}^{\delta_{ij}}t_i$ and $\gamma_j\cdot=
\sigma_j$ on $\g$, the $\sigma_j$ commuting of orders $m_j$.  Then
$A=k[t^{\pm1}]$ is reduced and finitely generated (H1); the action of
$\Gamma$ on $(k^\times)^n$ is free because $\prod_j\zeta_{m_j}^{a_j}=1$
componentwise forces $m_j\mid a_j$ (H2); $\g$ is simple (H3).
Theorem~\ref{thm:main} applies.

\subsection{Identification with the standard multiloop algebra}
Decompose $\g=\bigoplus_{\bar j\in\Gamma^\vee}\g_{\bar j}$ into
simultaneous eigenspaces of the commuting $\sigma_1,\dots,\sigma_n$, where
$\Gamma^\vee=\prod_j\Z/m_j$ and $\sigma_j|_{\g_{\bar k}}=
\zeta_{m_j}^{k_j}$.  As in the loop case, a Laurent monomial
$a\,t_1^{i_1}\cdots t_n^{i_n}$ ($a\in\g$) is $\Gamma$-invariant iff
$\sigma_j(a)=\zeta_{m_j}^{i_j}a$ for all $j$, i.e.\
$a\in\g_{(\bar i_1,\dots,\bar i_n)}$.  Hence
\[
\M^\Gamma=M(\g,\sigma_1,\dots,\sigma_n)
=\bigoplus_{i\in\Z^n}\g_{(\bar i_1,\dots,\bar i_n)}\,t_1^{i_1}\cdots
t_n^{i_n},
\]
the standard multiloop algebra.

\subsection{The variety automorphisms $\Aut((\Gm)^n,\Gamma)$}
$\Aut_{\mathrm{var}}((\Gm)^n)=(k^\times)^n\rtimes\GL_n(\Z)$:
a variety automorphism is a monomial map
$\varphi_{a,B}\colon t_i\mapsto a_i\prod_\ell t_\ell^{B_{i\ell}}$ with
$a\in(k^\times)^n$ and $B\in\GL_n(\Z)$.  (These are exactly the $k$-algebra
automorphisms of $A=k[t_1^{\pm1},\dots,t_n^{\pm1}]$: the group of units is
$A^\times=k^\times\times\{$Laurent monomials$\}\cong k^\times\times\Z^n$, an
algebra automorphism is determined by its effect on the generating units
$t_1,\dots,t_n$, hence has the displayed monomial form, and $B$ must be
invertible over $\Z$ since $\varphi^{-1}$ is again monomial.)

\emph{Precise normalization condition.}  The $\Gamma$-action factors through
a homomorphism into the torus: writing $\zeta_{m_j}$ for a fixed primitive
$m_j$-th root of unity, $\gamma_j$ acts as the translation by the point
$\rho_j=(1,\dots,\zeta_{m_j},\dots,1)\in(\Gm)^n$ (the $\zeta_{m_j}$ in slot
$j$), i.e.\ $\gamma_j\cdot t=\rho_j\!\cdot\!t$ coordinatewise.  Thus the
image of $\Gamma$ in $\Aut((\Gm)^n)$ is the finite \emph{translation}
subgroup $T_\Gamma=\{$translations by $\rho\in\mu\}$, where
$\mu=\prod_j\langle\zeta_{m_j}\rangle\cong\prod_j\Z/m_j\subseteq(\Gm)^n$ is
a finite subgroup of the torus.  For $\varphi=\varphi_{a,B}$ and a
translation $\tau_\rho\colon t\mapsto\rho t$ one computes
\[
\varphi\,\tau_\rho\,\varphi^{-1}(t)
=\varphi\bigl(\rho\,\varphi^{-1}(t)\bigr)
=\tau_{\,B\!\cdot\!\rho}(t),
\qquad\text{where }(B\!\cdot\!\rho)_i:=\prod_\ell\rho_\ell^{B_{i\ell}},
\]
because the linear scaling $a$ commutes with translations up to the same
shift and the monomial part sends a translation by $\rho$ to a translation
by $B\!\cdot\!\rho$ (the action of $B\in\GL_n(\Z)$ on the torus
$\mathrm{Hom}(\Z^n,\Gm)=(\Gm)^n$ via its natural action on the cocharacter
lattice $\Z^n$).  Hence $\varphi$ normalizes $T_\Gamma$ (equivalently the
$\Gamma$-action) iff
\begin{equation}\label{eq:Bnormalizes}
B\cdot\mu=\mu\quad\text{(as subgroups of }(\Gm)^n),
\end{equation}
i.e.\ iff the automorphism of the torus induced by $B$ preserves the finite
subgroup $\mu$.  Identifying $\mu\cong\Gamma$ via $\rho_j\leftrightarrow
\gamma_j$, condition \eqref{eq:Bnormalizes} means $B$ induces a group
automorphism $c_B\in\Aut(\Gamma)$ by $c_B(\gamma)\leftrightarrow B\!\cdot\!
(\text{the corresponding }\rho)$; explicitly, reducing $B\bmod m_j$ in each
relevant coordinate, $c_B$ is the automorphism of $\prod_j\Z/m_j$ induced by
$B$ acting on the quotient $\Z^n/(\text{lattice of characters trivial on }
\mu)\cong\prod_j\Z/m_j$.  Therefore
\[
\Aut((\Gm)^n,\Gamma)=
\bigl((k^\times)^n\rtimes\GL_n(\Z)\bigr)_{\Gamma}
:=\{(a,B):a\in(k^\times)^n,\ B\in\GL_n(\Z),\ B\cdot\mu=\mu\},
\]
with the surjection $(a,B)\mapsto c_B$ onto the subgroup of $\Aut(\Gamma)$
realized by $\GL_n(\Z)$-matrices preserving $\mu$.  (The scaling $a$ acts
trivially on $\Gamma$, so it lies in the kernel and contributes the
$(k^\times)^n$ of pure rescalings.)

\subsection{The gauge maps and the automorphism group}
By Theorem~\ref{thm:main} and Remark~\ref{rem:groupstr},
\[
\Aut\bigl(M(\g,\sigma_1,\dots,\sigma_n)\bigr)\cong
\Aut((\Gm)^n,\Gamma)\ \ltimes\ \mathcal G_\Gamma((\Gm)^n,\g),
\]
where the gauge group $\mathcal G_\Gamma((\Gm)^n,\g)$ consists of regular
maps $\Psi\colon(\Gm)^n\to G=\Aut(\g)$ satisfying, for $c=c_B\in\Aut(\Gamma)$,
\[
\Psi(\gamma x)=c(\gamma)\,\Psi(x)\,\gamma^{-1}\qquad(\gamma\in\Gamma),
\]
i.e.\ $\Psi$ is equivariant for the $\Gamma$-rotation on the source and the
$c$-twisted conjugation by $\langle\sigma_1,\dots,\sigma_n\rangle$ on the
target.  (As in the loop case, $\Psi$ is valued in all of $G$, not just in
$N_\Gamma(\g)$.)  As in \S\ref{sec:loop}, since $(\Gm)^n$ is connected and
$\mathrm{Out}(\g)$ is finite, each $\Psi$ factors as
$\Psi=\theta_0\cdot\Psi^\circ$ with $\theta_0$ a constant element of
$N_\Gamma(\g)$ (a representative of the component of $G$ containing the
connected image $\Psi((\Gm)^n)$) and $\Psi^\circ\colon(\Gm)^n\to
\mathrm{Inn}(\g)$ an inner gauge loop.  The identity
$\mathcal G^\circ=\mathrm{Inn}(M)$ for the multivariable inner gauge group
is proved exactly as \eqref{eq:innerloop}, with one honest caveat in the
generation step.  Part (i) (each generator $\exp(\mathrm{ad}\,\xi)$,
$\xi=v\,t^{i_1}_1\cdots t^{i_n}_n$ with $v\in\g_{(\bar i_1,\dots,\bar i_n)}$
ad-nilpotent, is an equivariant inner gauge monomial in $\mathrm{Inn}(M)$)
is the verbatim computation of the loop case and uses no external input; the
Steinberg torus reduction of Lemma~\ref{lem:steinberg} likewise carries over
verbatim with $R=A=k[t_1^{\pm1},\dots,t_n^{\pm1}]$.  Part (ii) uses the
multivariable analogue of Theorem~\ref{thm:looptorgen}: by the structure of
loop-reductive group schemes and Pianzola's $H^1$-vanishing for the iterated
Laurent ring $A^\Gamma$ (a Laurent polynomial ring in $n$ variables over
$k$, of finite global dimension), the twisted multiloop group
$\mathcal G^\circ=\underline{\bar G}(A^\Gamma)$ is generated by its
(relative) equivariant root subgroups $\exp(\mathrm{ad}\,\xi)$ together with
the equivariant torus elements of Lemma~\ref{lem:steinberg}, each of which
lies in $\mathrm{Inn}(M)$ by part (i).  \emph{Caveat (honest scope).}  For
$n\ge2$ the bare elementary-generation statement $\bar G(A)=E(A)\cdot\bar
T(A)$ over the higher-dimensional Laurent ring $A$ is more delicate than in
the one-variable case (it can fail for general split groups over multivariate
Laurent rings, reflecting nontrivial higher $K$-theory); what survives, and
what we use, is the generation of the \emph{twisted form's} group of
$A^\Gamma$-points by its relative root and torus subgroups, which is Pianzola's
loop-reductive generation theorem and is the precise multivariable
counterpart of Theorem~\ref{thm:looptorgen}.  We invoke it as an external
result and verify by hand, exactly as in the loop case, that each of its
generators lies in $\mathrm{Inn}(M)$.

\emph{Constant/diagram part and the $B$-compatibility.}  We make precise the
group $N_\Gamma(\g)$ that supplies the finite outer part.  Set
$\Sigma:=\langle\sigma_1,\dots,\sigma_n\rangle\subseteq\Aut(\g)$, a finite
abelian subgroup (the $\sigma_j$ commute, of orders $m_j$), and identify
$\Gamma\twoheadrightarrow\Sigma$ via $\gamma_j\mapsto\sigma_j$ (the
``equivariance data'' homomorphism).  When the $\sigma_j$ generate a copy
of $\prod_j\Z/m_j$ (the case where $\Gamma$ acts faithfully on $\g$) this
is an isomorphism $\Gamma\cong\Sigma$; in general $\Sigma$ is the quotient
of $\Gamma$ by the kernel of the action on $\g$, and the statements below
refer to $\Sigma=\mathrm{im}(\Gamma\to\Aut(\g))$ throughout (the gauge
equivariance \eqref{eq:equiv} only sees the images in $G$, so this causes no
change).  A constant gauge $\Psi\equiv\theta$
satisfies the twisted equivariance \eqref{eq:equiv} (for the variety part
$\varphi=\varphi_{a,B}$ with induced $c_B\in\Aut(\Gamma)$) iff
\[
\theta=c_B(\gamma)\,\theta\,\gamma^{-1}\quad\text{in }G
\quad\Longleftrightarrow\quad
\theta\,\sigma_\gamma\,\theta^{-1}=\sigma_{c_B(\gamma)}\quad(\gamma\in\Gamma),
\]
writing $\sigma_\gamma$ for the image of $\gamma$ in $\Sigma$.  Thus a
constant gauge $\theta$ is admissible \emph{together with} the variety part
$B$ iff conjugation by $\theta$ realizes on $\Sigma$ exactly the
automorphism $c_B$ of $\Gamma\cong\Sigma$.  In particular $\theta$ must
\emph{normalize} $\Sigma$, i.e.\ $\theta\in N_G(\Sigma)=N_\Gamma(\g)$, and
the resulting conjugation automorphism of $\Sigma$ must be realizable by some
$B\in\GL_n(\Z)$ preserving $\mu$ (Condition~\eqref{eq:Bnormalizes}).
Conversely, given $\theta\in N_\Gamma(\g)$, conjugation by $\theta$ is an
automorphism of the finite abelian group $\Sigma\cong\prod_j\Z/m_j$, hence
\emph{is} induced by some integer matrix $B\in\GL_n(\Z)$ preserving the
character data $\mu$ (every automorphism of $\prod_j\Z/m_j$ lifts to a
$\GL_n(\Z)$-matrix preserving $\mu$, since $\mathrm{Aut}(\prod_j\Z/m_j)$ is
the reduction of the stabilizer of $\mu$ in $\GL_n(\Z)$ acting on
$\Z^n\twoheadrightarrow\prod_j\Z/m_j$).  Pairing $\theta$ with this $B$ (and
any scaling $a$) gives an admissible constant gauge.  Therefore the
finite constant/diagram part of the gauge group is governed by
\[
\Aut(\g;\sigma_1,\dots,\sigma_n):=N_\Gamma(\g)=N_G(\Sigma)
=\{\theta\in\Aut(\g):\theta\,\Sigma\,\theta^{-1}=\Sigma\},
\]
and the compatibility ``$\theta$ permutes/powers the $\sigma_i$ as $B$ acts
on $\Gamma^\vee=\prod_j\Z/m_j$'' is exactly the statement that conjugation by
$\theta$ on $\Sigma$ equals $c_B$ on $\Gamma\cong\Sigma$, as just shown.  Assembling, every
automorphism of $M$ has the form $(\Phi f)(t)=\Psi(t)\,f(\varphi^{-1}t)$
with $\varphi$ monomial normalizing $\Gamma$ and $\Psi$ as above; the
group is
\[
\Aut(M)\cong
\bigl((k^\times)^n\rtimes\GL_n(\Z)\bigr)_\Gamma\ \ltimes\
\mathcal G_\Gamma((\Gm)^n,\g),
\qquad
\mathcal G_\Gamma((\Gm)^n,\g)=\mathrm{Inn}(M)\cdot R,
\]
with $R$ the finite outer/diagram part (the image of $N_\Gamma(\g)$ in
$\mathrm{Out}(\g)$).  This is the explicit description claimed in
Theorem~\ref{thm:multiloop}.

\section{Relation to the isomorphism classification}\label{sec:iso}

The same rigidity inputs (Lemmas~\ref{lem:centroid}--\ref{lem:fibrerigid})
classify \emph{isomorphisms} $\M(X,\g)^\Gamma\xrightarrow{\sim}
\M(X',\g')^{\Gamma'}$, not just automorphisms: replacing $\Phi$ by an
isomorphism between two equivariant map algebras, the centroid argument
gives a variety isomorphism $X/\Gamma\cong X'/\Gamma'$, the fibre argument
gives $\g\cong\g'$, and the equivariance \eqref{eq:equiv} matches the
$\Gamma$- and $\Gamma'$-structures.  Hence:

\begin{corollary}\label{cor:iso}
Under \textnormal{(H1)--(H3)} for both sides, two equivariant map algebras
$\M(X,\g)^\Gamma$ and $\M(X',\g')^{\Gamma'}$ are isomorphic as Lie algebras
if and only if there is a variety isomorphism $\varphi\colon X\to X'$
intertwining the $\Gamma$- and $\Gamma'$-actions up to an automorphism
$c\colon\Gamma\xrightarrow{\sim}\Gamma'$, together with a regular map
$\Psi\colon X\to\mathrm{Isom}_{\mathrm{Lie}}(\g',\g)$ satisfying the
twisted equivariance \eqref{eq:equiv}.  In particular:
\begin{itemize}
\item Two twisted loop algebras $L(\g,\sigma)$, $L(\g',\sigma')$ are
isomorphic iff $\g\cong\g'$ and $\langle\sigma\rangle$,
$\langle\sigma'\rangle$ are conjugate in $\Aut(\g)$ up to the order
$m\mapsto m$ data; equivalently iff the diagrams and the diagram-automorphism
classes agree (recovering the standard list of affine Kac--Moody
realizations $X_N^{(r)}$).
\item Two multiloop algebras $M(\g,\sigma_\bullet)$,
$M(\g',\sigma'_\bullet)$ are isomorphic iff $\g\cong\g'$ and the subgroups
$\langle\sigma_1,\dots,\sigma_n\rangle$,
$\langle\sigma'_1,\dots,\sigma'_n\rangle$ of $\Aut(\g)$ are conjugate
\emph{after} an admissible change of coordinates $B\in\GL_n(\Z)$ on the
grading lattice; this is the classification of multiloop algebras by the
$\GL_n(\Z)$-conjugacy class of the finite abelian subgroup
$\langle\sigma_\bullet\rangle\subseteq\Aut(\g)$.
\end{itemize}
\end{corollary}
\begin{proof}
Run the proof of Theorem~\ref{thm:main} with an isomorphism in place of an
automorphism: every step (perfectness, centroid, fibre, equivariance) used
only that the map is a Lie isomorphism, and produces the pair
$(\varphi,\Psi)$ relating the two structures.  Conversely such a pair
induces an isomorphism by the computation in the proof of
Theorem~\ref{thm:main}.  Specializing $X=X'=\Gm$ (resp.\ $(\Gm)^n$) and
reading off the constraints \eqref{eq:equiv} for the rotation actions gives
the conjugacy statements: for $\Gm$, $\varphi=t\mapsto at^{\pm1}$ forces
$\sigma'\sim\sigma^{\pm1}$ and $m=m'$; for $(\Gm)^n$, the matrix
$B\in\GL_n(\Z)$ implements the lattice change of coordinates, so the
grading groups $\langle\sigma_\bullet\rangle$ must be $\GL_n(\Z)$-conjugate
in $\Aut(\g)$.  The loop-algebra statement is the classical correspondence
between $\langle\sigma\rangle$-conjugacy classes (i.e.\ diagram
automorphisms of order $r$) and affine types $X_N^{(r)}$.
\end{proof}

\begin{remark}
The hypotheses (H1)--(H3) cover all the listed special cases.  When $\fa$
is reductive but not simple, or $\Gamma$ does not act freely on $X(k)$, the
same method applies after stratifying $X$ by orbit type and replacing the
single simple quotient in Lemma~\ref{lem:fibres} by a product of simple
quotients indexed by the orbit; the centroid then equals
$\OO(X/\Gamma)\otimes Z(\fa)$ and the gauge group is taken with values in
$\Aut(\fa)$.  The structure
``$\Aut=\,$(variety automorphisms normalizing $\Gamma$)$\,\ltimes\,$(twisted
equivariant gauge group)'' persists; this is the broad-generality answer.
\end{remark}

\end{document}
